\documentclass[a4paper,11pt]{article}
\usepackage[utf8]{inputenc}
\usepackage[T1]{fontenc}
\usepackage[italian]{babel}
\usepackage{lmodern}
\usepackage[margin=2cm]{geometry}
\usepackage{bpninja}
\usepackage{booktabs,tabularx,enumitem}
\usepackage[hidelinks]{hyperref}
\usepackage{xcolor}
\setlist{itemsep=1pt,topsep=3pt}
\newcommand{\nodo}[1]{\textsf{\textbf{#1}}}
\newcommand{\var}[1]{\texttt{#1}}
\newenvironment{ricetta}{\par\smallskip\noindent\textbf{Ingredienti.}\begin{itemize}}{\end{itemize}}
\newcommand{\ninja}[1]{\par\smallskip\noindent\colorbox{black!6}{\parbox{\dimexpr\linewidth-2\fboxsep}{\small\textbf{Mossa da ninja.} #1}}\par\smallskip}

\title{\textbf{Diventare ninja dei Blueprint}\\[4pt]\large Il platformer del corso, meccanica per meccanica, rifatto a nodi}
\author{Guida ai Blueprint del progetto \texttt{AIV\_UE5Intro}}
\date{}

\begin{document}
\maketitle
\tableofcontents
\newpage

\section{Come leggere questa guida}

Ogni meccanica del platformer (i tre livelli \texttt{L\_GD13}) è scritta nel progetto in C++, come componente o come actor. Qui la stessa logica è disegnata \emph{a nodi}, come la si costruisce nell'Event Graph di un Blueprint. Il personaggio si chiama \var{BP\_Eroe}: è un Blueprint figlio di \nodo{Character} (rinominatelo come volete). Ogni sezione ha:
\begin{itemize}
  \item gli \textbf{ingredienti}: in quale Blueprint si lavora, quali componenti e variabili servono;
  \item uno o più \textbf{grafici}, da ricopiare nodo per nodo;
  \item la \textbf{mossa da ninja}: il trucco che fa la differenza, o l'errore classico da evitare.
\end{itemize}
I nomi dei nodi sono quelli dell'editor in inglese (quelli che si cercano con il tasto destro nel grafo). I numeri sono gli stessi usati nei livelli: se li cambiate, cambiano le distanze che si riescono a saltare.

\subsection{Legenda}

\begin{center}
\begin{tabular}{ll@{\qquad}ll}
\toprule
\pin{exec} & esecuzione (filo nero) & \pin{bool} & Boolean \\
\pin{int} & Integer & \pin{float} & Float \\
\pin{vec} & Vector & \pin{rot} & Rotator \\
\pin{obj} & riferimento a un oggetto & \pin{class} & classe \\
\pin{str} & String & \pin{text} & Text \\
\pin{name} & Name & \pin{struct} & struttura (Hit Result\dots) \\
\pin{enum} & enum & \pin{del} & delegate / evento \\
\bottomrule
\end{tabular}
\end{center}

\begin{bpgraph}
\bpnode{a}{(0,0)}{bpev}{Evento}{-}{exec/}
\bpnode{b}{\dx{a}{0.8}{0}}{bpfn}{Funzione (impura)}{exec/,float/Ingresso}{exec/,float/Uscita}
\bpnode{c}{\dx{b}{0.8}{0}}{bppure}{Funzione pura}{float/A,float/B}{float/Risultato}
\bpnode{d}{\dx{c}{0.8}{0}}{bpflow}{Branch / Sequence}{exec/,bool/Condition}{exec/True,exec/False}
\bpvar{e}{(0,-1.8)}{float}{Variabile (Get)}
\bpnode{f}{\sotto{b}{0}{-0.5}}{bpcast}{Cast To \dots}{exec/,obj/Object}{exec/,exec/Cast Failed,obj/As \dots}
\bpnode{g}{\sotto{c}{0}{-0.5}}{bpcustom}{Evento custom}{-}{exec/}
\bpnode{h}{\sotto{d}{0}{-0.5}}{bptl}{Timeline}{exec/Play,exec/Reverse}{exec/Update,exec/Finished,float/Alpha}
\end{bpgraph}

Le testate colorate dicono che tipo di nodo è: \textcolor{bpev}{\textbf{rosso}} gli eventi, \textcolor{bpfn}{\textbf{blu}} le funzioni con il pin di esecuzione, \textcolor{bppure}{\textbf{verde}} le funzioni pure (senza esecuzione: si ricalcolano ogni volta che qualcuno legge il loro risultato), \textcolor{bpflow}{\textbf{grigio}} il controllo di flusso, \textcolor{bpcast}{\textbf{verde acqua}} i cast, \textcolor{bpcustom}{\textbf{viola}} eventi custom e dispatcher, \textcolor{bptl}{\textbf{arancio}} le timeline. Le capsule colorate sono le variabili lette con un \emph{Get}.

Quando un nodo verde ha per titolo una formula (per esempio \nodo{Z $-$ 380 $\cdot$ dt}) è un piccolo conto: potete farlo con i nodi aritmetici uno per uno, oppure con un unico nodo \nodo{Math Expression} (tasto destro $\rightarrow$ \emph{Add Math Expression}) scrivendo la formula così com'è. Le uscite degli \emph{Hit Result} si leggono con \nodo{Break Hit Result}; i vettori con \nodo{Break Vector} / \nodo{Make Vector}.

\subsection{Gli asset da creare prima}

\begin{center}
\begin{tabularx}{\linewidth}{lX}
\toprule
Asset & A cosa serve \\
\midrule
\var{BP\_Eroe} (Character) & il giocatore: movimento, abilità, torcia, vite, casse \\
\var{GI\_Platform} (Game Instance) & quello che sopravvive ai cambi di livello: vite, monete, abilità, record \\
\var{SG\_Platform} (Save Game) & il salvataggio su disco \\
\var{E\_Abilita} (Enumeration) & \var{DoppioSalto}, \var{WallRun}, \var{Jetpack}, \var{GeneratoreCasse} \\
\var{BPI\_Interagibile} (Blueprint Interface) & funzioni \var{Interagisci(Chi)} e \var{GetPrompt} \\
\var{IA\_Jump}, \var{IA\_Interact} & già nel progetto (Third Person) \\
\var{IA\_Zoom} & Input Action, Value Type \emph{Axis1D (float)} \\
\var{IA\_GeneraCassa} & Input Action, Value Type \emph{Digital (bool)} \\
\var{IMC\_Ninja} & Input Mapping Context: \var{IA\_Zoom} = Mouse Wheel Axis, \var{IA\_GeneraCassa} = F e Gamepad Face Button Left \\
\var{WBP\_HUD} (Widget Blueprint) & vite, monete, carburante, batteria, casse, timer \\
\var{GM\_Platform} (Game Mode) & Default Pawn Class = \var{BP\_Eroe}, HUD/Widget \\
\bottomrule
\end{tabularx}
\end{center}

In \nodo{Project Settings} $\rightarrow$ \nodo{Maps \& Modes} impostate \var{GI\_Platform} come Game Instance Class e \var{GM\_Platform} come Default GameMode.

\subsection{Il primo grafo: attaccare l'input}

Le azioni nuove (\var{IA\_Zoom}, \var{IA\_GeneraCassa}) stanno in un Mapping Context nostro, aggiunto \emph{sopra} a quello di base con una priorità più alta. In \var{BP\_Eroe}:

\begin{bpgraph}
\bpnode{bp}{(0,0)}{bpev}{Event BeginPlay}{-}{exec/}
\bpnode{pc}{\sotto{bp}{0.6}{-0.5}}{bppure}{Get Player Controller}{int/Player Index 0}{obj/Return Value}
\bpnode{ss}{\dx{pc}{0.6}{0}}{bppure}{Enhanced Input Local Player Subsystem}{obj/Player Controller}{obj/Return Value}
\bpnode{add}{\dx{ss}{0.8}{1.6}}{bpfn}{Add Mapping Context}{exec/,obj/Target,obj/Mapping Context: IMC\_Ninja,int/Priority 1}{exec/}
\bpexec{bp-out1}{add-in1}
\bpwire{obj}{pc-out1}{ss-in1}
\bpwire{obj}{ss-out1}{add-in2}
\end{bpgraph}

\ninja{Non toccate \var{IMC\_Default} del template: aggiungete un \emph{vostro} context. Se un giorno l'abilità non serve più, basta un \nodo{Remove Mapping Context} e i tasti tornano liberi.}

\section{La memoria del gioco: \texttt{GI\_Platform}}

Un livello si distrugge quando se ne apre un altro: tutto quello che vive dentro (anche \var{BP\_Eroe}) sparisce. La \nodo{Game Instance} invece nasce all'avvio del gioco e muore alla chiusura. È il posto giusto per vite, monete, abilità sbloccate e record. Nel C++ è \texttt{UPlatformerRunSubsystem}.

\begin{ricetta}
  \item Blueprint \var{GI\_Platform}, classe padre \nodo{GameInstance}.
  \item Variabili: \var{Vite} (Integer, 3), \var{MaxVite} (Integer, 3), \var{Monete} (Integer), \var{Abilita} (Set di \var{E\_Abilita}), \var{Salvataggio} (riferimento a \var{SG\_Platform}), \var{PrimoLivello} (Name, \var{L\_GD13\_Livello1}).
  \item \var{SG\_Platform} (classe padre \nodo{SaveGame}) con le stesse \var{Monete} e \var{Abilita}, più una Map \var{Record} (Name $\rightarrow$ Float).
  \item Event Dispatcher \var{OnAbilitaSbloccata} con un parametro \var{E\_Abilita}.
\end{ricetta}

\subsection{Caricare o creare il salvataggio}

\begin{bpgraph}
\bpnode{init}{(0,0)}{bpev}{Event Init}{-}{exec/}
\bpnode{ex}{\dx{init}{0.6}{0}}{bpfn}{Does Save Game Exist}{exec/,str/Slot Name "Platform",int/User Index 0}{exec/,bool/Return Value}
\bpnode[2.4]{br}{\dx{ex}{0.6}{0}}{bpflow}{Branch}{exec/,bool/Condition}{exec/True,exec/False}
\bpnode{load}{\dx{br}{0.6}{0}}{bpfn}{Load Game From Slot}{exec/,str/Slot Name,int/User Index}{exec/,obj/Return Value}
\bpnode{cast}{\dx{load}{0.6}{0}}{bpcast}{Cast To SG\_Platform}{exec/,obj/Object}{exec/,exec/Cast Failed,obj/As SG Platform}
\bpnode{set}{\dx{cast}{0.6}{0}}{bpvarh}{Set Salvataggio}{exec/,obj/Salvataggio}{exec/,obj/}
\bpnode{cre}{\sotto{load}{0}{-0.6}}{bpfn}{Create Save Game Object}{exec/,class/Class: SG\_Platform}{exec/,obj/Return Value}
\bpexec{init-out1}{ex-in1}
\bpexec{ex-out1}{br-in1}
\bpwire{bool}{ex-out2}{br-in2}
\bpexec{br-out1}{load-in1}
\bpexec{load-out1}{cast-in1}
\bpwire{obj}{load-out2}{cast-in2}
\bpexec{cast-out1}{set-in1}
\bpwire{obj}{cast-out3}{set-in2}
\bpexec{br-out2}{cre-in1}
\bpexec{cre-out1}{set-in1}
\bpwire{obj}{cre-out2}{set-in2}
\bpbox{All'avvio: se lo slot esiste lo carico, altrimenti ne creo uno nuovo}{(init-nw) (set-se) (cre-se)}
\end{bpgraph}

Subito dopo, copiate \var{Monete} e \var{Abilita} dal salvataggio nelle variabili della Game Instance (due \nodo{Set} che leggono \var{Salvataggio}).

\subsection{Sbloccare un'abilità}

Funzione \var{SbloccaAbilita(Quale)}: aggiunge l'enum al Set, salva subito e avvisa chi ascolta (HUD, personaggio).

\begin{bpgraph}
\bpnode{f}{(0,0)}{bpcustom}{SbloccaAbilita}{-}{exec/,enum/Quale}
\bpnode{add}{\dx{f}{0.6}{0}}{bpfn}{ADD (Set)}{exec/,wild/Target: Abilita,enum/New Item}{exec/}
\bpnode{cp}{\dx{add}{0.6}{0}}{bpvarh}{Set Abilita (SG)}{exec/,obj/Target: Salvataggio,wild/Abilita}{exec/}
\bpnode{sv}{\dx{cp}{0.6}{0}}{bpfn}{Save Game to Slot}{exec/,obj/Save Game Object,str/Slot Name "Platform",int/User Index 0}{exec/,bool/Return Value}
\bpnode{call}{\dx{sv}{0.6}{0}}{bpcustom}{Call OnAbilitaSbloccata}{exec/,enum/Quale}{exec/}
\bpexec{f-out1}{add-in1}\bpexec{add-out1}{cp-in1}\bpexec{cp-out1}{sv-in1}\bpexec{sv-out1}{call-in1}
\bpwire{enum}{f-out2}{add-in3}
\bpdraw[bpdataw, draw=bpc-enum] (f-out2) .. controls +(3,-2.6) and +(-3,-1.8) .. (call-in2);
\end{bpgraph}

Per sapere se il giocatore ha un'abilità basta una funzione \emph{pura} \var{HaAbilita(Quale)} $\rightarrow$ \nodo{CONTAINS (Set)}. Spuntate \emph{Pure} nei dettagli della funzione: così non ha pin di esecuzione e la potete collegare direttamente a un \nodo{Branch}.

\subsection{Perdere una vita e il Game Over}

\begin{bpgraph}
\bpnode{f}{(0,0)}{bpcustom}{Funzione PerdiVita}{-}{exec/}
\bpvar{v}{\sotto{f}{0.4}{-0.4}}{int}{Vite}
\bpnode[1.4]{dec}{\dx{v}{0.5}{0.1}}{bppure}{$-$}{int/,int/1}{int/}
\bpnode[2.4]{set}{\dx{f}{2.6}{0}}{bpvarh}{Set Vite}{exec/,int/Vite}{exec/,int/}
\bpnode[1.4]{gt}{\dx{set}{0.5}{-1.2}}{bppure}{$>$}{int/,int/0}{bool/}
\bpnode[2.4]{br}{\dx{set}{2.4}{0}}{bpflow}{Branch}{exec/,bool/Condition}{exec/True,exec/False}
\bpnode{rt}{\dx{br}{0.6}{0}}{bpflow}{Return Node}{exec/,bool/Continua: true}{-}
\bpnode{go}{\sotto{rt}{0}{-0.6}}{bpcustom}{GameOver}{exec/}{exec/}
\bpnode{rf}{\dx{go}{0.6}{0}}{bpflow}{Return Node}{exec/,bool/Continua: false}{-}
\bpwire{int}{v-out1}{dec-in1}
\bpwire{int}{dec-out1}{set-in2}
\bpexec{f-out1}{set-in1}
\bpwire{int}{set-out2}{gt-in1}
\bpexec{set-out1}{br-in1}
\bpwire{bool}{gt-out1}{br-in2}
\bpexec{br-out1}{rt-in1}
\bpexec{br-out2}{go-in1}
\bpexec{go-out1}{rf-in1}
\bpbox{PerdiVita (funzione con uscita Continua): True se si riparte dal punto sicuro}{(f-nw) (rf-se) (v-se)}
\end{bpgraph}

\var{GameOver} è un evento custom: mostra il messaggio, svuota le abilità (se volete la regola dura del C++, \var{bGameOverResetsAbilities}) e dopo due secondi riparte dal primo livello.

\begin{bpgraph}
\bpnode{go}{(0,0)}{bpcustom}{Event GameOver}{-}{exec/}
\bpnode{clr}{\dx{go}{0.6}{0}}{bpfn}{CLEAR (Set)}{exec/,wild/Target: Abilita}{exec/}
\bpnode[2.4]{d}{\dx{clr}{0.6}{0}}{bpflow}{Delay}{exec/,float/Duration 2.0}{exec/Completed}
\bpnode{sv}{\dx{d}{0.6}{0}}{bpvarh}{Set Vite}{exec/,int/= MaxVite}{exec/}
\bpnode{op}{\dx{sv}{0.6}{0}}{bpfn}{Open Level (by Name)}{exec/,name/Level Name: PrimoLivello}{exec/}
\bpexec{go-out1}{clr-in1}\bpexec{clr-out1}{d-in1}\bpexec{d-out1}{sv-in1}\bpexec{sv-out1}{op-in1}
\end{bpgraph}

\subsection{Cambiare livello}

\var{VaiAlLivello(Nome)}: \nodo{Set Vite} = \var{MaxVite}, poi \nodo{Open Level (by Name)}. Così ogni livello parte con tre vite, come nel C++ (\texttt{GoToLevel}). Da qualsiasi Blueprint si arriva qui con \nodo{Get Game Instance} $\rightarrow$ \nodo{Cast To GI\_Platform}.

\ninja{Il cast alla Game Instance fatelo \emph{una volta} nel BeginPlay di \var{BP\_Eroe} e salvate il risultato in una variabile \var{GI}. Castare a ogni Tick funziona, ma è come chiedere il nome a qualcuno ogni volta che lo incontri.}

\section{Il movimento: salti, wall run e jetpack}

Nel C++ è tutto in \texttt{UPlatformerAbilitiesComponent} e \texttt{UPlatformerJetpackComponent}. In Blueprint conviene creare un \nodo{Actor Component} \var{AC\_Movimento} da aggiungere a \var{BP\_Eroe}: il personaggio resta pulito e la stessa logica si può dare a un altro Character. Dentro il componente il personaggio è \nodo{Get Owner} $\rightarrow$ \nodo{Cast To Character}, salvato in BeginPlay nelle variabili \var{Eroe} e \var{Move} (il suo \emph{Character Movement}).

\subsection{La taratura}

Prima di scrivere nodi si toccano i numeri, in \var{BP\_Eroe} $\rightarrow$ componente \emph{Character Movement}:

\begin{center}
\begin{tabular}{lrl}
\toprule
Proprietà & Valore & Effetto \\
\midrule
Jump Z Velocity & 650 & salto alto circa 1,6 m \\
Air Control & 0,4 & si corregge la traiettoria in aria, ma non troppo \\
Max Walk Speed & 600 & corsa \\
Rotation Rate (Z) & 720 & il personaggio si gira subito \\
Orient Rotation To Movement & sì & guarda dove va (terza persona) \\
\bottomrule
\end{tabular}
\end{center}

\subsection{Coyote time e jump buffer}

Due trucchi che il giocatore non vede ma sente. \emph{Coyote time}: per 0,15\,s dopo essere scesi da un bordo si può ancora saltare. \emph{Jump buffer}: se si preme salto 0,15\,s prima di toccare terra, il salto parte appena si atterra.
Variabili: \var{UltimoAppoggio}, \var{UltimaRichiesta} (Float, secondi di gioco), \var{SaltiInAria} (Integer).

\begin{bpgraph}
\bpnode{tick}{(0,0)}{bpev}{Event Tick}{-}{exec/,float/Delta Seconds}
\bpnode{gr}{\sotto{tick}{0.3}{-0.4}}{bppure}{Is Moving On Ground}{obj/Target: Move}{bool/Return Value}
\bpnode[2.4]{br}{\dx{tick}{2.2}{0}}{bpflow}{Branch}{exec/,bool/Condition}{exec/True,exec/False}
\bpnode{gt}{\sotto{br}{0}{-0.4}}{bppure}{Get Game Time in Seconds}{-}{float/Return Value}
\bpnode{set}{\dx{br}{1.6}{0}}{bpvarh}{Set UltimoAppoggio}{exec/,float/UltimoAppoggio}{exec/}
\bpexec{tick-out1}{br-in1}\bpwire{bool}{gr-out1}{br-in2}
\bpexec{br-out1}{set-in1}\bpwire{float}{gt-out1}{set-in2}
\bpbox{Ogni frame: ricordo l'ultimo istante in cui avevo i piedi per terra}{(tick-nw) (set-se) (gt-se) (gr-se)}
\end{bpgraph}

Il salto non passa più direttamente a \nodo{Jump}: passa da una funzione \var{ProvaASaltare}, che restituisce \emph{True} se un salto è partito.

\begin{bpgraph}
\bpnode{ia}{(0,0)}{bpev}{EnhancedInputAction IA\_Jump}{-}{exec/Triggered,exec/Started,exec/Completed}
\bpnode{try}{\dx{ia}{0.6}{-0.4}}{bpfn}{ProvaASaltare}{exec/}{exec/,bool/Saltato}
\bpnode[2.4]{br}{\dx{try}{0.6}{0}}{bpflow}{Branch}{exec/,bool/Condition}{exec/True,exec/False}
\bpnode{gt}{\sotto{try}{0}{-0.5}}{bppure}{Get Game Time in Seconds}{-}{float/Return Value}
\bpnode{set}{\dx{br}{0.6}{-0.4}}{bpvarh}{Set UltimaRichiesta}{exec/,float/UltimaRichiesta}{exec/}
\bpexec{ia-out2}{try-in1}\bpexec{try-out1}{br-in1}\bpwire{bool}{try-out2}{br-in2}
\bpexec{br-out2}{set-in1}\bpwire{float}{gt-out1}{set-in2}
\bpbox{Se non riesco a saltare adesso, mi segno l'ora della richiesta (jump buffer)}{(ia-nw) (set-se) (gt-se)}
\end{bpgraph}

Dentro \var{ProvaASaltare} si provano, in ordine, i tre modi di saltare. Il primo che riesce esce con \emph{True}.

\begin{bpgraph}
\bpnode{f}{(0,0)}{bpcustom}{ProvaASaltare}{-}{exec/}
\bpnode{cj}{\sotto{f}{0}{-0.4}}{bppure}{Can Jump}{obj/Target: Eroe}{bool/Return Value}
\bpnode{coy}{\sotto{cj}{0}{-0.3}}{bppure}{GameTime $-$ UltimoAppoggio $<$ 0.15}{-}{bool/Coyote}
\maxx{m1}{cj,coy}
\bpnode{or}{\dxy{m1}{cj}{0.6}{-0.2}}{bppure}{OR Boolean}{bool/A,bool/B}{bool/}
\bpnode[2.4]{br}{\dxy{or}{f}{0.6}{0}}{bpflow}{Branch}{exec/,bool/Condition}{exec/True,exec/False}
\bpnode{jmp}{\dx{br}{0.6}{0}}{bpfn}{Launch Character}{exec/,obj/Target: Eroe,vec/Velocity (0 0 650),bool/XY Override false,bool/Z Override true}{exec/}
\bpnode{rt}{\dx{jmp}{0.6}{0}}{bpflow}{Return Node}{exec/,bool/Saltato: true}{-}
\bpnode{dj}{\sotto{jmp}{0}{-0.5}}{bpfn}{ProvaDoppioSalto}{exec/}{exec/,bool/Saltato}
\bpnode{rt2}{\dx{dj}{0.6}{0}}{bpflow}{Return Node}{exec/,bool/Saltato}{-}
\bpexec{f-out1}{br-in1}
\bpwire{bool}{cj-out1}{or-in1}\bpwire{bool}{coy-out1}{or-in2}\bpwire{bool}{or-out1}{br-in2}
\bpexec{br-out1}{jmp-in1}\bpexec{jmp-out1}{rt-in1}
\bpexec{br-out2}{dj-in1}\bpexec{dj-out1}{rt2-in1}\bpwire{bool}{dj-out2}{rt2-in2}
\end{bpgraph}

Perché \nodo{Launch Character} e non \nodo{Jump}? Perché durante il coyote time il personaggio è già \emph{Falling} e \nodo{Jump} lo rifiuterebbe. Con \emph{Z Override} la velocità verticale viene sostituita: il salto è identico a quello da terra.

Il buffer si chiude all'atterraggio:

\begin{bpgraph}
\bpnode{l}{(0,0)}{bpev}{Event On Landed}{-}{exec/,struct/Hit}
\bpnode{s0}{\dx{l}{0.6}{0}}{bpvarh}{Set SaltiInAria}{exec/,int/0}{exec/}
\bpnode{buf}{\sotto{l}{0}{-0.6}}{bppure}{GameTime $-$ UltimaRichiesta $<$ 0.15}{-}{bool/}
\maxx{m1}{s0,buf}
\bpnode[2.4]{br}{\dxy{m1}{s0}{0.6}{0}}{bpflow}{Branch}{exec/,bool/Condition}{exec/True,exec/False}
\bpnode{j}{\dx{br}{0.6}{0}}{bpfn}{Jump}{exec/,obj/Target: Eroe}{exec/}
\bpnode{clr}{\dx{j}{0.6}{0}}{bpvarh}{Set UltimaRichiesta}{exec/,float/$-$100}{exec/}
\bpexec{l-out1}{s0-in1}\bpexec{s0-out1}{br-in1}\bpwire{bool}{buf-out1}{br-in2}
\bpexec{br-out1}{j-in1}\bpexec{j-out1}{clr-in1}
\end{bpgraph}
(\nodo{Event On Landed} esiste solo nel Character: se lavorate nel componente, in BeginPlay fate \nodo{Bind Event to Landed} sull'\var{Eroe}.)

\ninja{Azzerate \var{UltimaRichiesta} dopo averla usata, altrimenti un atterraggio subito dopo fa partire un secondo salto ``fantasma''.}

\subsection{Il doppio salto controllabile}

Il secondo salto in aria non è un semplice \nodo{Jump}: prende la direzione in cui il giocatore spinge lo stick \emph{in quel momento}. Così si può cambiare direzione a mezz'aria, come nel C++ (\texttt{DoubleJumpSteering = 1}). Valori: 1 salto in aria, Z = 600, velocità orizzontale almeno 350.

\begin{bpgraph}
\bpnode{f}{(0,0)}{bpcustom}{ProvaDoppioSalto}{-}{exec/}
\bpnode{ha}{\sotto{f}{0}{-0.4}}{bppure}{HaAbilita}{enum/DoppioSalto}{bool/}
\bpnode{lt}{\sotto{ha}{0}{-0.3}}{bppure}{SaltiInAria $<$ 1}{-}{bool/}
\maxx{m1}{ha,lt}
\bpnode{and}{\dxy{m1}{ha}{0.6}{0}}{bppure}{AND Boolean}{bool/,bool/}{bool/}
\bpnode[2.4]{br}{\dxy{and}{f}{0.6}{0}}{bpflow}{Branch}{exec/,bool/Condition}{exec/True,exec/False}
\bpnode{rf}{\sotto{lt}{0}{-0.5}}{bpflow}{Return Node}{exec/,bool/Saltato: false}{-}
\bpnode{in}{\dx{br}{0.6}{-0.2}}{bppure}{Get Last Movement Input Vector}{obj/Target: Eroe}{vec/Return Value}
\bpnode{vel}{\sotto{in}{0}{-0.3}}{bppure}{Max(VLungo(Velocity XY) , 350)}{-}{float/Velocita}
\bpnode{mk}{\dx{in}{0.6}{0}}{bppure}{Dir $\cdot$ Velocita + (0 0 600)}{vec/Dir,float/Velocita}{vec/}
\bpexec{f-out1}{br-in1}\bpwire{bool}{ha-out1}{and-in1}\bpwire{bool}{lt-out1}{and-in2}\bpwire{bool}{and-out1}{br-in2}
\bpdraw[bpexecw] (br-out2) .. controls +(0.8,0) and +(-4,0) .. (rf-in1);
\bpwire{vec}{in-out1}{mk-in1}\bpwire{float}{vel-out1}{mk-in2}
\end{bpgraph}

\begin{bpgraph}
\bpnode{x}{(0,0)}{bpflow}{(dal True)}{-}{exec/}
\bpnode{lc}{\dx{x}{0.6}{0}}{bpfn}{Launch Character}{exec/,obj/Target: Eroe,vec/Launch Velocity,bool/XY Override true,bool/Z Override true}{exec/}
\bpnode{inc}{\dx{lc}{0.6}{0}}{bpfn}{Increment Int}{exec/,int/SaltiInAria}{exec/}
\bpnode{call}{\dx{inc}{0.6}{0}}{bpcustom}{Call OnDoppioSalto}{exec/}{exec/}
\bpnode{rt}{\dx{call}{0.6}{0}}{bpflow}{Return Node}{exec/,bool/Saltato: true}{-}
\bpexec{x-out1}{lc-in1}\bpexec{lc-out1}{inc-in1}\bpexec{inc-out1}{call-in1}\bpexec{call-out1}{rt-in1}
\end{bpgraph}

Se il giocatore non tocca lo stick, \nodo{Get Last Movement Input Vector} è zero e il salto va dritto in su: è il comportamento giusto. Nel C++ durante il doppio salto l'\emph{Air Control} sale a 0,8 finché non si atterra (\nodo{Set Air Control} sul \var{Move}, e in \nodo{On Landed} si rimette 0,4).

\subsection{La corsa sul muro}

Ogni frame, se si è in aria, due \nodo{Line Trace By Channel} cercano un muro a destra e a sinistra (80\,cm). Il muro vale solo se ha il tag \var{WallRun} e se è davvero verticale (la Z della normale è quasi zero). Variabili: \var{InWallRun} (Bool), \var{NormaleMuro} (Vector), \var{TempoSulMuro} (Float), \var{UltimoMuro} (Actor).

\begin{bpgraph}
\bpnode{tick}{(0,0)}{bpev}{Event Tick}{-}{exec/,float/Delta Seconds}
\bpnode{fall}{\sotto{tick}{0}{-0.4}}{bppure}{Is Falling}{obj/Target: Move}{bool/}
\bpnode[2.4]{br}{\dx{tick}{0.6}{0}}{bpflow}{Branch}{exec/,bool/Condition}{exec/True,exec/False}
\bpnode{loc}{\sotto{br}{0}{-0.9}}{bppure}{Get Actor Location}{obj/Target: Eroe}{vec/}
\bpnode{rv}{\sotto{loc}{0}{-0.3}}{bppure}{Get Actor Right Vector}{obj/Target: Eroe}{vec/}
\maxx{m1}{loc,rv}
\bpnode{end}{\dxy{m1}{loc}{0.6}{-0.2}}{bppure}{Loc + Right $\cdot$ 80}{vec/Loc,vec/Right}{vec/}
\bpnode{lt}{\dxy{end}{tick}{0.6}{0}}{bpfn}{Line Trace By Channel}{exec/,vec/Start,vec/End,enum/Visibility,bool/Ignore Self}{exec/,struct/Out Hit,bool/Return Value}
\bpexec{tick-out1}{br-in1}\bpwire{bool}{fall-out1}{br-in2}
\bpexec{br-out1}{lt-in1}\bpwire{vec}{loc-out1}{lt-in2}\bpwire{vec}{loc-out1}{end-in1}\bpwire{vec}{rv-out1}{end-in2}\bpwire{vec}{end-out1}{lt-in3}
\bpbox{Cerco il muro a destra (a sinistra: lo stesso, con Right $\cdot$ $-$80)}{(tick-nw) (lt-se) (rv-se)}
\end{bpgraph}

\begin{bpgraph}
\bpnode{bh}{(0,0)}{bpstruct}{Break Hit Result}{struct/Out Hit}{bool/Blocking Hit,vec/Impact Normal,obj/Hit Actor}
\bpnode{nz}{\dx{bh}{0.8}{0.6}}{bppure}{Abs(Normal.Z) $<$ 0.3}{vec/Normal}{bool/}
\bpnode{tag}{\sotto{nz}{0}{-0.3}}{bppure}{Actor Has Tag}{obj/Target,name/Tag: WallRun}{bool/}
\bpnode{sp}{\sotto{tag}{0}{-0.3}}{bppure}{VLungo(Velocity XY) $>$ 200}{-}{bool/}
\maxx{m1}{nz,tag,sp}
\bpnode{and}{\dxy{m1}{tag}{0.6}{0}}{bppure}{AND (4 ingressi)}{bool/Hit,bool/Verticale,bool/Tag,bool/Veloce}{bool/}
\bpnode[2.4]{br}{\dx{and}{0.6}{0.4}}{bpflow}{Branch}{exec/,bool/Condition}{exec/True,exec/False}
\bpnode{st}{\dx{br}{0.6}{0}}{bpcustom}{IniziaWallRun}{exec/,vec/Normale,obj/Muro}{exec/}
\bpwire{bool}{bh-out1}{and-in1}\bpwire{vec}{bh-out2}{nz-in1}\bpwire{obj}{bh-out3}{tag-in1}
\bpwire{bool}{nz-out1}{and-in2}\bpwire{bool}{tag-out1}{and-in3}\bpwire{bool}{sp-out1}{and-in4}
\bpwire{bool}{and-out1}{br-in2}\bpexec{br-out1}{st-in1}
\end{bpgraph}
(Per aggiungere ingressi a un \nodo{AND} si clicca \emph{Add pin +}.)

Durante la corsa la gravità la gestiamo noi: \nodo{Set Gravity Scale} a 0 in \var{IniziaWallRun}, poi a ogni Tick si impone la velocità lungo il muro e una discesa lenta.

\begin{bpgraph}
\bpnode{t}{(0,0)}{bpflow}{Tick, se InWallRun}{-}{exec/,float/dt}
\bpnode{dir}{\sotto{t}{0}{-0.4}}{bppure}{Cross Product}{vec/A: NormaleMuro,vec/B: (0 0 1)}{vec/}
\bpnode{sg}{\dx{dir}{0.6}{0}}{bppure}{Dir $\cdot$ Sign(Dot(Dir;Velocity))}{vec/Dir}{vec/Lungo}
\bpnode{z}{\sotto{sg}{0}{-0.3}}{bppure}{Max(Vel.Z $-$ 380 $\cdot$ dt ; $-$900)}{float/dt}{float/Z}
\bpnode{mk}{\dx{sg}{0.6}{0}}{bppure}{Lungo $\cdot$ 650 + (0 0 Z)}{vec/Lungo,float/Z}{vec/}
\maxx{m1}{mk,z}
\bpnode{sv}{\dxy{m1}{t}{0.6}{0}}{bpvarh}{Set Velocity}{exec/,obj/Target: Move,vec/Velocity}{exec/}
\bpnode{add}{\dx{sv}{0.6}{0}}{bpfn}{TempoSulMuro += dt}{exec/}{exec/}
\bpnode[2.4]{br}{\dx{add}{0.6}{0}}{bpflow}{Branch}{exec/,bool/$>$ 1.5 s}{exec/True,exec/False}
\bpnode{fine}{\sotto{br}{0.4}{-0.4}}{bpcustom}{FineWallRun}{exec/}{exec/}
\bpexec{t-out1}{sv-in1}\bpwire{vec}{dir-out1}{sg-in1}\bpwire{vec}{sg-out1}{mk-in1}\bpwire{float}{z-out1}{mk-in2}\bpwire{vec}{mk-out1}{sv-in3}
\bpexec{sv-out1}{add-in1}\bpexec{add-out1}{br-in1}\bpdraw[bpexecw] (br-out1) .. controls +(0.8,0) and +(-0.8,0) .. (fine-in1);
\end{bpgraph}

\var{FineWallRun} rimette \nodo{Set Gravity Scale} a 1, \var{InWallRun} a false e salva il muro in \var{UltimoMuro}: per 0,45\,s quel muro non si può riprendere (altrimenti il personaggio ci resta incollato). Toccare un muro ridà anche il salto in aria (\var{SaltiInAria} = 0).

\subsection{Il salto dal muro}

In \var{ProvaASaltare}, \emph{prima} di tutto il resto, si controlla \var{InWallRun}. Se è vero, il salto spinge via dal muro e in alto:

\begin{bpgraph}
\bpnode[2.4]{br}{(0,0)}{bpflow}{Branch}{exec/,bool/InWallRun}{exec/True,exec/False}
\bpnode{v}{\sotto{br}{0}{-0.4}}{bppure}{Normale $\cdot$ 600 + (0 0 700) + Lungo $\cdot$ 650 $\cdot$ 0.5}{-}{vec/}
\bpnode{fine}{\dxy{v}{br}{0.6}{0}}{bpcustom}{FineWallRun}{exec/}{exec/}
\bpnode{lc}{\dx{fine}{0.6}{0}}{bpfn}{Launch Character}{exec/,obj/Eroe,vec/Launch Velocity,bool/XY Override true,bool/Z Override true}{exec/}
\bpnode{rt}{\dx{lc}{0.6}{0}}{bpflow}{Return Node}{exec/,bool/Saltato: true}{-}
\bpexec{br-out1}{fine-in1}\bpexec{fine-out1}{lc-in1}\bpexec{lc-out1}{rt-in1}\bpwire{vec}{v-out1}{lc-in3}
\end{bpgraph}

\subsection{Il jetpack: tenere premuto salto}

Il jetpack non ha un tasto suo: si \emph{tiene premuto} salto. Dopo 0,2\,s di pressione in aria la spinta si accende. La spinta non aggiunge forza: porta la velocità verticale verso +260\,cm/s con \nodo{FInterp To Constant} (accelerazione 2400). Così il jetpack frena una caduta e sale piano, ma non si trasforma in un razzo. Variabili: \var{Premuto} (Bool), \var{TempoPremuto}, \var{Carburante} (2,2\,s), \var{CarburanteMax}.

\begin{bpgraph}
\bpnode{ia}{(0,0)}{bpev}{EnhancedInputAction IA\_Jump}{-}{exec/Triggered,exec/Started,exec/Completed}
\bpnode{s1}{\dx{ia}{0.6}{0.2}}{bpvarh}{Set Premuto}{exec/,bool/true}{exec/}
\bpnode{s2}{\dx{s1}{0.6}{0}}{bpvarh}{Set TempoPremuto}{exec/,float/0}{exec/}
\bpnode{s3}{\sotto{s1}{0}{-0.4}}{bpvarh}{Set Premuto}{exec/,bool/false}{exec/}
\bpexec{ia-out2}{s1-in1}\bpexec{s1-out1}{s2-in1}\bpexec{ia-out3}{s3-in1}
\end{bpgraph}

\begin{bpgraph}
\bpnode{t}{(0,0)}{bpev}{Event Tick}{-}{exec/,float/Delta Seconds}
\bpnode{c}{\sotto{t}{0}{-0.4}}{bppure}{Premuto AND Is Falling AND HaAbilita(Jetpack)}{-}{bool/}
\bpnode[2.4]{br}{\dxy{c}{t}{0.6}{0}}{bpflow}{Branch}{exec/,bool/Condition}{exec/True,exec/False}
\bpnode{add}{\dx{br}{0.6}{0}}{bpfn}{TempoPremuto += dt}{exec/}{exec/}
\bpnode{c2}{\sotto{add}{0}{-0.4}}{bppure}{TempoPremuto $>$ 0.2 AND Carburante $>$ 0}{-}{bool/}
\bpnode[2.4]{br2}{\dxy{c2}{add}{0.6}{0}}{bpflow}{Branch}{exec/,bool/Condition}{exec/True,exec/False}
\bpexec{t-out1}{br-in1}\bpwire{bool}{c-out1}{br-in2}\bpexec{br-out1}{add-in1}\bpexec{add-out1}{br2-in1}\bpwire{bool}{c2-out1}{br2-in2}
\end{bpgraph}

\begin{bpgraph}
\bpnode{x}{(0,0)}{bpflow}{(dal True)}{-}{exec/}
\bpnode{vz}{\sotto{x}{0}{-0.4}}{bppure}{FInterp To Constant}{float/Current: Velocity.Z,float/Target 260,float/Delta Time,float/Interp Speed 2400}{float/}
\bpnode{mk}{\dx{vz}{0.6}{0}}{bppure}{Make Vector (Vel.X; Vel.Y; Z)}{float/Z}{vec/}
\bpnode{sv}{\dxy{mk}{x}{0.6}{0}}{bpvarh}{Set Velocity}{exec/,obj/Target: Move,vec/Velocity}{exec/}
\bpnode{fuel}{\dx{sv}{0.6}{0}}{bpvarh}{Set Carburante}{exec/,float/Carburante $-$ dt}{exec/}
\bpnode{fx}{\dx{fuel}{0.6}{0}}{bpfn}{Set Visibility (Fiamma)}{exec/,bool/New Visibility: true}{exec/}
\bpexec{x-out1}{sv-in1}\bpwire{float}{vz-out1}{mk-in1}\bpwire{vec}{mk-out1}{sv-in3}\bpexec{sv-out1}{fuel-in1}\bpexec{fuel-out1}{fx-in1}
\end{bpgraph}

Sul ramo \emph{False} la fiamma si spegne. Il carburante si ricarica solo a terra: nel ramo \emph{True} del Tick del coyote time aggiungete \nodo{Set Carburante} = \nodo{Min(Carburante + 1.5 $\cdot$ dt ; CarburanteMax)}. La fiamma è un cono (\var{SM\_Cone}) con un materiale emissivo arancione e una \nodo{Point Light} figlia, attaccati alla schiena di \var{BP\_Eroe}.

\ninja{Usate \nodo{FInterp To Constant} e non \nodo{FInterp To}: la versione ``Constant'' cambia la velocità sempre dello stesso passo (2400 al secondo), quella normale rallenta man mano che si avvicina al bersaglio e la spinta sembrerebbe molle.}

\section{La torcia e la prima persona con la rotella}

La rotella del mouse accorcia e allunga il braccio della telecamera (\emph{Spring Arm}). Quando il braccio scende sotto i 40\,cm la telecamera entra nella testa: il personaggio sparisce, la visuale diventa in prima persona e si accende la torcia. Nel C++ è \texttt{UPlatformerTorchComponent}.

\begin{ricetta}
  \item In \var{BP\_Eroe}: il \var{CameraBoom} (Spring Arm) e la \var{FollowCamera} ci sono già nel template.
  \item Sotto la \var{FollowCamera}: \var{Torcia} (Static Mesh, un cilindro scuro ruotato di 90\textdegree, posizione $(30, 18, -18)$) e \var{LuceTorcia} (\nodo{Spot Light}: Intensity 5000, Attenuation Radius 2500, Inner/Outer Cone 18\textdegree/32\textdegree, colore caldo). Entrambi \emph{Visible} = false.
  \item Variabili: \var{BraccioVoluto} (Float, 400), \var{BraccioMax} (400), \var{PrimaPersona} (Bool), \var{Batteria} (Float, 25), \var{BatteriaMax} (25).
  \item \var{IA\_Zoom} (Axis1D) mappata su \emph{Mouse Wheel Axis} in \var{IMC\_Ninja}.
\end{ricetta}

\subsection{La rotella cambia il braccio}

Ogni scatto della rotella vale $\pm 1$: lo moltiplichiamo per 60\,cm e teniamo il risultato fra 0 e \var{BraccioMax}.

\begin{bpgraph}
\bpnode{ia}{(0,0)}{bpev}{EnhancedInputAction IA\_Zoom}{-}{exec/Triggered,exec/Started,exec/Completed,float/Action Value}
\bpvar{bv}{\sotto{ia}{0}{-0.4}}{float}{BraccioVoluto}
\maxx{m1}{ia,bv}
\bpnode{calc}{\dxy{m1}{bv}{0.6}{0.4}}{bppure}{BraccioVoluto $-$ Value $\cdot$ 60}{float/BraccioVoluto,float/Value}{float/}
\bpnode{cl}{\dx{calc}{0.6}{0}}{bppure}{Clamp (float)}{float/Value,float/Min 0,float/Max: BraccioMax}{float/}
\bpnode{set}{\dxy{cl}{ia}{0.6}{0}}{bpvarh}{Set BraccioVoluto}{exec/,float/}{exec/}
\bpexec{ia-out1}{set-in1}
\bpwire{float}{bv-out1}{calc-in1}\bpwire{float}{ia-out4}{calc-in2}\bpwire{float}{calc-out1}{cl-in1}\bpwire{float}{cl-out1}{set-in2}
\end{bpgraph}
(Se la rotella zooma al contrario di come vi aspettate, cambiate il segno: $+$ al posto di $-$.)

\subsection{Il braccio scivola verso il valore voluto}

Nel Tick il braccio non salta al valore nuovo: ci arriva morbido con \nodo{FInterp To}. Poi si controlla la soglia della prima persona.

\begin{bpgraph}
\bpnode{t}{(0,0)}{bpev}{Event Tick}{-}{exec/,float/Delta Seconds}
\bpnode{fi}{\sotto{t}{0}{-0.4}}{bppure}{FInterp To}{float/Current: Target Arm Length,float/Target: BraccioVoluto,float/Delta Time,float/Interp Speed 10}{float/}
\bpnode{set}{\dxy{fi}{t}{0.6}{0}}{bpvarh}{Set Target Arm Length}{exec/,obj/Target: CameraBoom,float/}{exec/}
\bpnode{lt}{\sotto{set}{0}{-0.4}}{bppure}{Target Arm Length $<$ 40}{-}{bool/}
\maxx{m1}{set,lt}
\bpnode{ne}{\dxy{m1}{t}{0.6}{0}}{bppure}{NOT EQUAL (Boolean)}{bool/A,bool/B: PrimaPersona}{bool/}
\bpnode[2.4]{br}{\dx{ne}{0.6}{0}}{bpflow}{Branch}{exec/,bool/Condition}{exec/True,exec/False}
\bpnode{call}{\dx{br}{0.6}{0}}{bpcustom}{ImpostaPrimaPersona}{exec/,bool/Attiva}{exec/}
\bpexec{t-out1}{set-in1}\bpwire{float}{fi-out1}{set-in3}
\bpwire{bool}{lt-out1}{ne-in1}\bpwire{bool}{ne-out1}{br-in2}
\bpdraw[bpexecw] (set-out1) .. controls +(2,0) and +(-2,0) .. (br-in1);
\bpexec{br-out1}{call-in1}
\bpdraw[bpdataw, draw=bpc-bool] (lt-out1) .. controls +(3,-0.5) and +(-1.5,-1) .. (call-in2);
\end{bpgraph}

Il controllo con \nodo{NOT EQUAL} serve a chiamare \var{ImpostaPrimaPersona} solo quando lo stato \emph{cambia}, non a ogni frame.

\subsection{Entrare nella testa}

\begin{bpgraph}
\bpnode{f}{(0,0)}{bpcustom}{ImpostaPrimaPersona}{-}{exec/,bool/Attiva}
\bpnode{s0}{\dx{f}{0.6}{0}}{bpvarh}{Set PrimaPersona}{exec/,bool/Attiva}{exec/}
\bpnode{ns}{\dx{s0}{0.6}{0}}{bpfn}{Set Owner No See}{exec/,obj/Target: Mesh,bool/New Owner No See: Attiva}{exec/}
\bpnode{off}{\dx{ns}{0.6}{0}}{bpvarh}{Set Target Offset}{exec/,obj/CameraBoom,vec/Select: (0 0 65) o (0 0 0)}{exec/}
\bpexec{f-out1}{s0-in1}\bpexec{s0-out1}{ns-in1}\bpexec{ns-out1}{off-in1}
\end{bpgraph}

\begin{bpgraph}
\bpnode{x}{(0,0)}{bpflow}{(continua)}{-}{exec/}
\bpnode{yaw}{\dx{x}{0.6}{0}}{bpvarh}{Set Use Controller Rotation Yaw}{exec/,obj/Target: self,bool/Attiva}{exec/}
\bpnode{orient}{\dx{yaw}{0.6}{0}}{bpvarh}{Set Orient Rotation To Movement}{exec/,obj/Target: Move,bool/NOT Attiva}{exec/}
\bpnode{v1}{\sotto{yaw}{0}{-0.5}}{bpfn}{Set Visibility}{exec/,obj/Target: Torcia,bool/Attiva}{exec/}
\bpnode{v2}{\dx{v1}{0.6}{0}}{bpfn}{Set Visibility}{exec/,obj/Target: LuceTorcia,bool/Attiva AND Batteria $>$ 0}{exec/}
\bpexec{x-out1}{yaw-in1}\bpexec{yaw-out1}{orient-in1}
\bpdraw[bpexecw] (orient-out1) .. controls +(1,0) and +(-1.5,0) .. ++(0,-0.9) .. controls +(-3,0) and +(-1,0) .. (v1-in1);
\bpexec{v1-out1}{v2-in1}
\end{bpgraph}

Perché funziona: \nodo{Set Owner No See} nasconde la mesh \emph{solo} alla telecamera del proprietario (le ombre restano). \emph{Target Offset} alza la telecamera all'altezza degli occhi (65\,cm sopra il centro della capsula). \emph{Use Controller Rotation Yaw} fa girare il corpo con il mouse, come in ogni FPS; in terza persona si torna a \emph{Orient Rotation To Movement}.

\subsection{La batteria}

\begin{bpgraph}
\bpnode{t}{(0,0)}{bpflow}{Tick (dopo lo zoom)}{-}{exec/,float/dt}
\bpnode[2.4]{br}{\dx{t}{0.6}{0}}{bpflow}{Branch}{exec/,bool/PrimaPersona}{exec/True,exec/False}
\bpnode{sub}{\dx{br}{0.6}{0}}{bpvarh}{Set Batteria}{exec/,float/Max(Batteria $-$ dt ; 0)}{exec/}
\bpnode{rec}{\sotto{sub}{0}{-0.6}}{bpvarh}{Set Batteria}{exec/,float/Min(Batteria + 0.25 $\cdot$ dt ; Max)}{exec/}
\bpnode{pct}{\dx{sub}{0.6}{-0.3}}{bppure}{Batteria / BatteriaMax}{-}{float/Carica}
\bpnode{fl}{\sotto{pct}{0}{-0.3}}{bppure}{Carica $<$ 0.2 ? RandomFloatInRange(0.3;1) : 1}{float/Carica}{float/Fattore}
\maxx{m1}{pct,fl}
\bpnode{int}{\dxy{m1}{sub}{0.6}{0}}{bpfn}{Set Intensity}{exec/,obj/Target: LuceTorcia,float/5000 $\cdot$ Fattore}{exec/}
\bpexec{t-out1}{br-in1}\bpexec{br-out1}{sub-in1}\bpexec{br-out2}{rec-in1}
\bpexec{sub-out1}{int-in1}\bpwire{float}{pct-out1}{fl-in1}\bpwire{float}{fl-out1}{int-in3}
\end{bpgraph}

Sotto il 20\% la luce tremola (un'intensità casuale ogni frame), e a zero si spegne (\nodo{Set Visibility} false). Il pickup di batteria (capitolo~\ref{sec:pickup}) aggiunge 15\,s.

\ninja{Il clamp sul braccio è la vera ``leva'': non servono due telecamere né un cambio di Pawn. Una sola Spring Arm con lunghezza variabile copre terza persona, inquadratura ravvicinata e prima persona.}

\section{Tre vite e il posto sicuro}

Cadere nel vuoto non fa ricominciare il livello: si perde una vita e si torna all'ultimo punto \emph{sicuro} su cui si stava in piedi. Finite le tre vite, si riparte dal livello 1. Nel C++ è \texttt{UPlatformerSafeSpotComponent} più \texttt{AFallZone}.

\begin{ricetta}
  \item Actor Component \var{AC\_PostoSicuro} su \var{BP\_Eroe}. Variabili: \var{PostoSicuro} (Vector), \var{HoPosto} (Bool), \var{Margine} (Float, 50), \var{CadutaMax} (Float, 2500).
  \item Un tag \var{NoSafeSpot} da mettere sugli actor su cui non si deve ricomparire (piattaforme mobili, che crollano, porte).
  \item Blueprint \var{BP\_ZonaCaduta}: un \nodo{Box Collision} grande sotto il livello.
\end{ricetta}

\subsection{Ricordare dove si stava bene}

Non serve il Tick: un timer ogni 0,2\,s basta e costa molto meno.

\begin{bpgraph}
\bpnode{bp}{(0,0)}{bpev}{Event BeginPlay}{-}{exec/}
\bpnode{tm}{\dx{bp}{0.6}{0}}{bpfn}{Set Timer by Event}{exec/,del/Event,float/Time 0.2,bool/Looping true}{exec/,struct/Return Value}
\bpnode{ev}{\sotto{bp}{0}{-1.7}}{bpcustom}{Custom Event: Campiona}{-}{exec/,del/}
\bpnode{gr}{\sotto{ev}{0}{-0.4}}{bppure}{Is Moving On Ground}{obj/Move}{bool/}
\maxx{m1}{ev,gr}
\bpnode[2.4]{br}{\dxy{m1}{ev}{0.6}{0}}{bpflow}{Branch}{exec/,bool/Condition}{exec/True,exec/False}
\bpnode{ok}{\dx{br}{0.6}{0}}{bpfn}{PavimentoSicuro?}{exec/}{exec/,bool/Sicuro}
\bpnode[2.4]{br2}{\dx{ok}{0.6}{0}}{bpflow}{Branch}{exec/,bool/Condition}{exec/True,exec/False}
\bpnode{set}{\dx{br2}{0.6}{0}}{bpvarh}{Set PostoSicuro}{exec/,vec/Get Actor Location}{exec/}
\bpexec{bp-out1}{tm-in1}
\bpdraw[bpdataw, draw=bpc-del] (ev-out2) .. controls +(1,0) and +(-1,0) .. (tm-in2);
\bpexec{ev-out1}{br-in1}\bpwire{bool}{gr-out1}{br-in2}\bpexec{br-out1}{ok-in1}\bpexec{ok-out1}{br2-in1}\bpwire{bool}{ok-out2}{br2-in2}\bpexec{br2-out1}{set-in1}
\end{bpgraph}
(Dopo \var{Set PostoSicuro} mettete anche \var{HoPosto} = true.)

\subsection{Che cosa è un pavimento sicuro}

Non basta essere a terra: se si è sul bordo, ricomparire lì vuol dire ricadere subito. La funzione \var{PavimentoSicuro?} spara quattro raggi verso il basso, spostati di 50\,cm avanti, indietro, a destra e a sinistra: devono colpire \emph{tutti} un pavimento buono.

\begin{bpgraph}
\bpnode{f}{(0,0)}{bpcustom}{PavimentoSicuro?}{-}{exec/}
\bpnode{arr}{\sotto{f}{0}{-0.4}}{bppure}{Make Array}{vec/(50 0 0),vec/($-$50 0 0),vec/(0 50 0),vec/(0 $-$50 0)}{wild/Array}
\bpnode{fe}{\dxy{arr}{f}{0.6}{0}}{bpflow}{For Each Loop}{exec/Exec,wild/Array}{exec/Loop Body,wild/Array Element,int/Array Index,exec/Completed}
\bpnode{lt}{\dx{fe}{0.6}{0}}{bpfn}{Line Trace By Channel}{exec/,vec/Start: Loc + Offset,vec/End: Start $-$ (0 0 150),enum/Visibility,bool/Ignore Self}{exec/,struct/Out Hit,bool/Return Value}
\bpexec{f-out1}{fe-in1}\bpwire{wild}{arr-out1}{fe-in2}\bpexec{fe-out1}{lt-in1}
\end{bpgraph}

\begin{bpgraph}
\bpnode{bh}{(0,0)}{bpstruct}{Break Hit Result}{struct/Out Hit}{bool/Blocking Hit,obj/Hit Actor,obj/Hit Component}
\bpnode{c1}{\dx{bh}{0.6}{0.6}}{bppure}{NOT Actor Has Tag (NoSafeSpot)}{obj/Hit Actor}{bool/}
\bpnode{c2}{\sotto{c1}{0}{-0.3}}{bppure}{NOT Is Valid (Cast To BP\_PiattaformaMobile)}{obj/Hit Actor}{bool/}
\bpnode{c3}{\sotto{c2}{0}{-0.3}}{bppure}{Component Velocity (VLungo) $<$ 1}{obj/Hit Component}{bool/}
\maxx{m1}{c1,c2,c3}
\bpnode{and}{\dxy{m1}{c1}{0.6}{-0.3}}{bppure}{AND}{bool/Blocking Hit,bool/,bool/,bool/}{bool/Buono}
\bpnode[2.4]{br}{\dxy{and}{bh}{0.6}{0}}{bpflow}{Branch}{exec/,bool/Buono}{exec/True,exec/False}
\bpnode{rf}{\dx{br}{0.6}{0}}{bpflow}{Return Node}{exec/,bool/Sicuro: false}{-}
\bpwire{bool}{bh-out1}{and-in1}\bpwire{obj}{bh-out2}{c1-in1}\bpwire{obj}{bh-out2}{c2-in1}\bpwire{obj}{bh-out3}{c3-in1}
\bpwire{bool}{c1-out1}{and-in2}\bpwire{bool}{c2-out1}{and-in3}\bpwire{bool}{c3-out1}{and-in4}\bpwire{bool}{and-out1}{br-in2}
\bpdraw[bpexecw] (br-out2) .. controls +(0.5,0) and +(-0.5,0) .. (rf-in1);
\end{bpgraph}
Se un raggio fallisce si esce subito con \emph{False}; se il \nodo{For Each} arriva a \emph{Completed} si esce con \emph{True}.

\subsection{La zona di caduta}

\begin{bpgraph}
\bpnode{ov}{(0,0)}{bpev}{On Component Begin Overlap (Box)}{-}{exec/,obj/Other Actor}
\bpnode{cast}{\dx{ov}{0.6}{0}}{bpcast}{Cast To BP\_Eroe}{exec/,obj/Object}{exec/,exec/Cast Failed,obj/As BP Eroe}
\bpnode{gc}{\dx{cast}{0.6}{0}}{bpfn}{GestisciCaduta}{exec/,obj/Target: AC\_PostoSicuro}{exec/}
\bpexec{ov-out1}{cast-in1}\bpwire{obj}{ov-out2}{cast-in2}\bpexec{cast-out1}{gc-in1}\bpwire{obj}{cast-out3}{gc-in2}
\end{bpgraph}

\begin{bpgraph}
\bpnode{f}{(0,0)}{bpcustom}{GestisciCaduta}{-}{exec/}
\bpnode{pv}{\dx{f}{0.6}{0}}{bpfn}{PerdiVita}{exec/,obj/Target: GI}{exec/,bool/Continua}
\bpnode{c}{\sotto{pv}{0}{-0.4}}{bppure}{Continua AND HoPosto}{bool/}{bool/}
\maxx{m1}{pv,c}
\bpnode[2.4]{br}{\dxy{m1}{f}{0.6}{0}}{bpflow}{Branch}{exec/,bool/Condition}{exec/True,exec/False}
\bpnode{sl}{\dx{br}{0.6}{0}}{bpfn}{Set Actor Location}{exec/,obj/Target: Eroe,vec/PostoSicuro + (0 0 20),bool/Teleport true}{exec/}
\bpnode{sv}{\dx{sl}{0.6}{0}}{bpvarh}{Set Velocity}{exec/,obj/Move,vec/(0 0 0)}{exec/}
\bpexec{f-out1}{pv-in1}\bpexec{pv-out1}{br-in1}\bpwire{bool}{pv-out2}{c-in1}\bpwire{bool}{c-out1}{br-in2}\bpexec{br-out1}{sl-in1}\bpexec{sl-out1}{sv-in1}
\end{bpgraph}

Sul ramo \emph{False} non serve fare niente: se le vite sono finite, \var{GameOver} della Game Instance sta già riaprendo il primo livello. Come rete di sicurezza, nel timer \var{Campiona} controllate anche la quota: se l'eroe è sceso più di \var{CadutaMax} (2500\,cm) sotto \var{PostoSicuro}, chiamate \var{GestisciCaduta} lo stesso (serve nei punti dove avete dimenticato la zona di caduta).

\ninja{Il tag \var{NoSafeSpot} è il modo più economico di dire ``qui no'' senza toccare il codice: lo si mette con due clic nel pannello Details $\rightarrow$ Actor $\rightarrow$ Tags. Fatelo anche sulle casse generate, altrimenti si ricompare sopra una cassa che poi sparisce.}

\section{Il generatore di casse}

Premendo \textbf{F} (o il tasto sinistro del pad) compare una cassa davanti all'eroe. Le cariche sono tre; quando finiscono si aspettano 5\,s e tornano tutte. Le casse si impilano, così si sale su muri che col salto non si passano (il muro di 4,6\,m del livello 1). Nel C++: \texttt{UPlatformerCrateSpawnerComponent} e \texttt{ASpawnedCrate}.

\begin{ricetta}
  \item Blueprint \var{BP\_Cassa} (Actor): un cubo 80\,cm (scala 0,8) con collisione \emph{BlockAll}, tag \var{NoSafeSpot}, \emph{Initial Life Span} = 25.
  \item In \var{BP\_Eroe} (o in un componente \var{AC\_Casse}): \var{Cariche} (Integer, 3), \var{CaricheMax} (3), \var{InRicarica} (Bool), \var{CasseVive} (Array di \var{BP\_Cassa}).
  \item \var{IA\_GeneraCassa} (Digital) mappata su F e \emph{Gamepad Face Button Left}.
\end{ricetta}

\subsection{Dove mettere la cassa}

Il punto davanti all'eroe è a 130\,cm. Da lì un raggio scende fino a 3\,m: se colpisce qualcosa (pavimento o un'altra cassa) la cassa si appoggia sopra; altrimenti compare in aria e cade (ha la fisica spenta, quindi resta lì: è voluto, si può usare come gradino a mezz'aria).

\begin{bpgraph}
\bpnode{ia}{(0,0)}{bpev}{EnhancedInputAction IA\_GeneraCassa}{-}{exec/Triggered,exec/Started,exec/Completed}
\bpnode{c}{\sotto{ia}{0}{-0.4}}{bppure}{Cariche $>$ 0 AND HaAbilita(GeneratoreCasse)}{-}{bool/}
\maxx{m1}{ia,c}
\bpnode[2.4]{br}{\dxy{m1}{ia}{0.6}{0}}{bpflow}{Branch}{exec/,bool/Condition}{exec/True,exec/False}
\bpnode{p}{\sotto{br}{0}{-0.8}}{bppure}{Loc + Forward $\cdot$ 130 + (0 0 150)}{-}{vec/Alto}
\bpnode{lt}{\dxy{p}{br}{0.6}{0}}{bpfn}{Line Trace By Channel}{exec/,vec/Start: Alto,vec/End: Alto $-$ (0 0 300),enum/Visibility}{exec/,struct/Out Hit,bool/Return Value}
\bpexec{ia-out2}{br-in1}\bpwire{bool}{c-out1}{br-in2}\bpexec{br-out1}{lt-in1}\bpwire{vec}{p-out1}{lt-in2}
\end{bpgraph}

\begin{bpgraph}
\bpnode{bh}{(0,0)}{bpstruct}{Break Hit Result}{struct/Out Hit}{bool/Blocking Hit,vec/Impact Point}
\bpnode{sel}{\dx{bh}{0.6}{0}}{bppure}{Select Vector}{vec/A: Impact Point + (0 0 40),vec/B: Alto,bool/Pick A: Blocking Hit}{vec/}
\bpnode{mt}{\dx{sel}{0.6}{0}}{bppure}{Make Transform}{vec/Location,rot/Rotation: Yaw dell'eroe,vec/Scale (1 1 1)}{struct/}
\bpnode{sp}{\dx{mt}{0.6}{0.4}}{bpfn}{SpawnActor BP\_Cassa}{exec/,class/Class,struct/Spawn Transform,enum/Collision: Adjust If Possible}{exec/,obj/Return Value}
\bpwire{vec}{bh-out2}{sel-in1}\bpwire{bool}{bh-out1}{sel-in3}\bpwire{vec}{sel-out1}{mt-in1}\bpwire{struct}{mt-out1}{sp-in3}
\end{bpgraph}

Il \emph{Collision Handling Override} su \emph{Try To Adjust Location, But Always Spawn} è importante: se la cassa tocca appena l'eroe, viene spostata invece di non comparire.

\subsection{Contare le cariche e ricaricare}

\begin{bpgraph}
\bpnode{x}{(0,0)}{bpflow}{(dopo lo Spawn)}{-}{exec/,obj/Cassa}
\bpnode{add}{\dx{x}{0.6}{0}}{bpfn}{ADD (Array)}{exec/,wild/CasseVive,obj/Cassa}{exec/}
\bpnode{dec}{\dx{add}{0.6}{0}}{bpfn}{Decrement Int}{exec/,int/Cariche}{exec/,int/}
\bpnode{c}{\sotto{dec}{0}{-0.4}}{bppure}{Cariche == 0 AND NOT InRicarica}{int/}{bool/}
\maxx{m1}{dec,c}
\bpnode[2.4]{br}{\dxy{m1}{x}{0.6}{0}}{bpflow}{Branch}{exec/,bool/Condition}{exec/True,exec/False}
\bpnode{s}{\dx{br}{0.6}{0}}{bpvarh}{Set InRicarica}{exec/,bool/true}{exec/}
\bpnode{tm}{\dx{s}{0.6}{0}}{bpfn}{Set Timer by Function Name}{exec/,str/Function Name: Ricarica,float/Time 5.0,bool/Looping false}{exec/}
\bpexec{x-out1}{add-in1}\bpwire{obj}{x-out2}{add-in3}\bpexec{add-out1}{dec-in1}\bpwire{int}{dec-out2}{c-in1}
\bpexec{dec-out1}{br-in1}\bpwire{bool}{c-out1}{br-in2}\bpexec{br-out1}{s-in1}\bpexec{s-out1}{tm-in1}
\end{bpgraph}

La funzione \var{Ricarica} rimette \var{Cariche} = \var{CaricheMax} e \var{InRicarica} = false. Per non riempire il livello, se \var{CasseVive} ha più di 6 elementi distruggete la più vecchia (\nodo{GET} indice 0 $\rightarrow$ \nodo{Destroy Actor} $\rightarrow$ \nodo{Remove Index} 0). Le casse che muoiono da sole dopo 25\,s lasciano un riferimento non valido nell'array: ripulitelo con \nodo{Is Valid} prima di usarlo.

\ninja{Il timer si avvia solo quando le cariche arrivano a zero, non a ogni cassa. È la regola del brief (``3 volte, poi 5 secondi''), ed è più leggibile per il giocatore: l'HUD mostra ``Ricarica\dots'' e un numero che scende (\nodo{Get Timer Remaining Time by Function Name}).}

\section{Gli actor della palestra}

Tre actor semplici, ognuno con un'idea da imparare: leggere la distanza dal giocatore, usare un timer che si ripete, decidere a caso all'inizio. Nel C++ sono in \texttt{PlatformerGym.h}.

\subsection{L'actor che cresce con la distanza}

Il faro del livello 1 è piccolo quando ci si è vicino e diventa grande da lontano: si vede sempre, e indica la meta. Variabili: \var{DistMin} (300), \var{DistMax} (3000), \var{ScalaMin} (1), \var{ScalaMax} (4).

\begin{bpgraph}
\bpnode{t}{(0,0)}{bpev}{Event Tick}{-}{exec/,float/Delta Seconds}
\bpnode{pc}{\sotto{t}{0}{-0.4}}{bppure}{Get Player Character}{int/0}{obj/}
\bpnode{d}{\dxy{pc}{pc}{0.6}{0}}{bppure}{Get Distance To}{obj/Target: self,obj/Other Actor}{float/}
\bpnode{mr}{\dx{d}{0.6}{0}}{bppure}{Map Range Clamped}{float/Value,float/In Range A: DistMin,float/In Range B: DistMax,float/Out Range A: ScalaMin,float/Out Range B: ScalaMax}{float/}
\bpnode{fi}{\dx{mr}{0.6}{0}}{bppure}{FInterp To}{float/Current: Scala,float/Target,float/Delta Time,float/Interp Speed 4}{float/}
\bpnode{set}{\dxy{fi}{t}{0.6}{0}}{bpfn}{Set Actor Scale 3D}{exec/,obj/Target: self,vec/New Scale (Scala Scala Scala)}{exec/}
\bpexec{t-out1}{set-in1}\bpwire{obj}{pc-out1}{d-in2}\bpwire{float}{d-out1}{mr-in1}\bpwire{float}{mr-out1}{fi-in2}\bpwire{float}{fi-out1}{set-in3}
\end{bpgraph}

\nodo{Map Range Clamped} è il nodo chiave di tutta la guida: trasforma un intervallo in un altro (300--3000\,cm diventa 1--4) e non esce mai dai limiti. Il \nodo{FInterp To} evita gli scatti quando il giocatore si teletrasporta.

\subsection{L'actor che cambia modello ogni secondo}

Il totem e la piattaforma ``mutante'' passano da cubo a cilindro a sfera a cono. Variabili: \var{Modelli} (Array di Static Mesh), \var{Materiali} (Array di Material), \var{Indice} (Integer).

\begin{bpgraph}
\bpnode{bp}{(0,0)}{bpev}{Event BeginPlay}{-}{exec/}
\bpnode{tm}{\dx{bp}{0.6}{0}}{bpfn}{Set Timer by Event}{exec/,del/Event,float/Time 1.0,bool/Looping true}{exec/}
\bpnode{ev}{\sotto{bp}{0}{-1.7}}{bpcustom}{Custom Event: Cambia}{-}{exec/,del/}
\bpnode{idx}{\sotto{ev}{0}{-0.4}}{bppure}{(Indice + 1) \% LENGTH(Modelli)}{-}{int/}
\maxx{m1}{ev,idx}
\bpnode{s}{\dxy{m1}{ev}{0.6}{0}}{bpvarh}{Set Indice}{exec/,int/}{exec/,int/}
\bpnode{g1}{\sotto{s}{0}{-0.4}}{bppure}{GET (Modelli)}{int/Indice}{obj/}
\maxx{m2}{s,g1}
\bpnode{sm}{\dxy{m2}{ev}{0.6}{0}}{bpfn}{Set Static Mesh}{exec/,obj/Target: Mesh,obj/New Mesh}{exec/}
\bpnode{mat}{\dx{sm}{0.6}{0}}{bpfn}{Set Material}{exec/,obj/Target: Mesh,int/Index 0,obj/GET (Materiali; Indice)}{exec/}
\bpexec{bp-out1}{tm-in1}\bpdraw[bpdataw, draw=bpc-del] (ev-out2) .. controls +(1,0) and +(-1,0) .. (tm-in2);
\bpexec{ev-out1}{s-in1}\bpwire{int}{idx-out1}{s-in2}\bpwire{int}{s-out2}{g1-in1}\bpexec{s-out1}{sm-in1}\bpwire{obj}{g1-out1}{sm-in3}\bpexec{sm-out1}{mat-in1}
\end{bpgraph}

Il \% (modulo) fa tornare l'indice a 0 dopo l'ultimo modello. Se l'actor è una piattaforma, cambiate anche la collisione: con \nodo{Set Static Mesh} la collisione segue la nuova mesh, quindi sulla sfera si scivola giù e sul cono si sta solo in punta. È lì il divertimento.

\subsection{L'actor che c'è solo nel 35\% delle partite}

Il ponte del livello 1 a volte c'è e a volte no: se non c'è, si passa dai sassi sotto. Si decide una volta sola, all'inizio.

\begin{bpgraph}
\bpnode{bp}{(0,0)}{bpev}{Event BeginPlay}{-}{exec/}
\bpnode{rb}{\sotto{bp}{0}{-0.4}}{bppure}{Random Bool with Weight}{float/Weight 0.35}{bool/}
\maxx{m1}{bp,rb}
\bpnode[2.4]{br}{\dxy{m1}{bp}{0.6}{0}}{bpflow}{Branch}{exec/,bool/Condition}{exec/True,exec/False}
\bpnode{msg}{\dx{br}{0.6}{0}}{bpfn}{Print String / mostra cartello}{exec/,str/"Il ponte c'è!"}{exec/}
\bpnode{hd}{\sotto{msg}{0}{-0.4}}{bpfn}{Set Actor Hidden In Game}{exec/,bool/true}{exec/}
\bpnode{col}{\dx{hd}{0.6}{0}}{bpfn}{Set Actor Enable Collision}{exec/,bool/false}{exec/}
\bpexec{bp-out1}{br-in1}\bpwire{bool}{rb-out1}{br-in2}\bpexec{br-out1}{msg-in1}\bpexec{br-out2}{hd-in1}\bpexec{hd-out1}{col-in1}
\end{bpgraph}

Si poteva anche usare \nodo{Destroy Actor}, ma nascondere e togliere la collisione permette di mostrare il ponte più tardi (per esempio con una leva). Nel C++ il messaggio arriva con un piccolo ritardo, quando l'HUD è già pronto: in Blueprint basta un \nodo{Delay} di 0,5\,s prima del messaggio.

\ninja{\nodo{Random Bool with Weight} con 0,35 dà \emph{True} il 35\% delle volte. Per provare tutti e due i casi senza riavviare cento volte, aggiungete una variabile \var{Forza} (enum: Caso, SempreSì, SempreNo) modificabile nell'editor.}

\section{Gli ostacoli}

\subsection{La torretta}

La torretta vede il giocatore entro 18\,m, gira la testa verso di lui con una velocità limitata e spara ogni 1,4\,s solo se la mira è buona (entro 12\textdegree) e nessun muro è in mezzo. Nel gym del livello 2 il colpo è mortale: fa ricominciare il livello. Nel C++: \texttt{ATurret} e \texttt{APlatformerProjectile}.

\begin{ricetta}
  \item \var{BP\_Torretta}: \var{Base} (cilindro), \var{Testa} (Scene Component) con figli \var{TestaMesh}, \var{Canna} e \var{Bocca} (Arrow o Scene, sulla punta della canna), \var{Rilevamento} (Sphere Collision, raggio 1800).
  \item Variabili: \var{Bersaglio} (Actor), \var{CadenzaPronta} (Bool), \var{RiavviaLivello} (Bool, \emph{Instance Editable}).
  \item \var{BP\_Proiettile}: \var{Sfera} (Sphere Collision, raggio 12, root), mesh, \nodo{Projectile Movement} (Initial/Max Speed 1500, Gravity Scale 0), \emph{Initial Life Span} 4.
\end{ricetta}

\begin{bpgraph}
\bpnode{ov}{(0,0)}{bpev}{On Component Begin Overlap (Rilevamento)}{-}{exec/,obj/Other Actor}
\bpnode{c}{\dx{ov}{0.6}{0}}{bpcast}{Cast To BP\_Eroe}{exec/,obj/Object}{exec/,exec/Cast Failed,obj/As BP Eroe}
\bpnode{s}{\dx{c}{0.6}{0}}{bpvarh}{Set Bersaglio}{exec/,obj/}{exec/}
\bpnode{ov2}{\sotto{ov}{0}{-0.5}}{bpev}{On Component End Overlap (Rilevamento)}{-}{exec/,obj/Other Actor}
\bpnode{s2}{\dxy{ov2}{ov2}{0.6}{0}}{bpvarh}{Set Bersaglio}{exec/,obj/None}{exec/}
\bpexec{ov-out1}{c-in1}\bpwire{obj}{ov-out2}{c-in2}\bpexec{c-out1}{s-in1}\bpwire{obj}{c-out3}{s-in2}\bpexec{ov2-out1}{s2-in1}
\end{bpgraph}

\begin{bpgraph}
\bpnode{t}{(0,0)}{bpev}{Event Tick}{-}{exec/,float/Delta Seconds}
\bpnode{iv}{\dx{t}{0.6}{0}}{bpflow}{Is Valid}{exec/,obj/Bersaglio}{exec/Is Valid,exec/Is Not Valid}
\bpnode{flr}{\sotto{t}{0}{-0.6}}{bppure}{Find Look at Rotation}{vec/Start: Testa World Location,vec/Target: Bersaglio Loc + (0 0 $-$60)}{rot/}
\bpnode{ri}{\dx{flr}{0.6}{0}}{bppure}{RInterp To}{rot/Current: Testa World Rotation,rot/Target,float/Delta Time,float/Interp Speed 3}{rot/}
\bpnode{cl}{\dx{ri}{0.6}{0}}{bppure}{Clamp Pitch ($-$40 ; 30)}{rot/}{rot/}
\maxx{m1}{iv,cl}
\bpnode{set}{\dxy{m1}{t}{0.6}{0}}{bpfn}{Set World Rotation}{exec/,obj/Target: Testa,rot/New Rotation}{exec/}
\bpexec{t-out1}{iv-in1}\bpexec{iv-out1}{set-in1}\bpwire{rot}{flr-out1}{ri-in2}\bpwire{rot}{ri-out1}{cl-in1}\bpwire{rot}{cl-out1}{set-in3}
\end{bpgraph}

Il \nodo{RInterp To} con velocità 3 è ciò che rende la torretta \emph{battibile}: se correte di lato, la testa resta indietro. Il $-60$ sulla quota mira al petto e non alla testa (la capsula è alta 180\,cm, il centro è a 90).

\begin{bpgraph}
\bpnode{x}{(0,0)}{bpflow}{(dopo Set World Rotation)}{-}{exec/}
\bpnode{ang}{\sotto{x}{0}{-0.4}}{bppure}{Angolo(Forward Testa ; Dir. Bersaglio) $<$ 12}{-}{bool/Mira}
\bpnode{vis}{\sotto{ang}{0}{-0.3}}{bppure}{Line Trace Bocca$\rightarrow$Bersaglio: Hit Actor == Bersaglio}{-}{bool/Vede}
\bpnode{and}{\dxy{vis}{ang}{0.6}{0}}{bppure}{AND}{bool/Mira,bool/Vede,bool/CadenzaPronta}{bool/}
\bpnode[2.4]{br}{\dxy{and}{x}{0.6}{0}}{bpflow}{Branch}{exec/,bool/Condition}{exec/True,exec/False}
\bpnode{fire}{\dx{br}{0.6}{0}}{bpcustom}{Spara}{exec/}{exec/}
\bpexec{x-out1}{br-in1}\bpwire{bool}{ang-out1}{and-in1}\bpwire{bool}{vis-out1}{and-in2}\bpwire{bool}{and-out1}{br-in2}\bpexec{br-out1}{fire-in1}
\end{bpgraph}
(L'angolo si calcola con \nodo{Dot Product} dei due vettori normalizzati e \nodo{Acos (Degrees)}.)

\begin{bpgraph}
\bpnode{f}{(0,0)}{bpcustom}{Spara}{-}{exec/}
\bpnode{s0}{\dx{f}{0.6}{0}}{bpvarh}{Set CadenzaPronta}{exec/,bool/false}{exec/}
\bpnode{tr}{\sotto{s0}{0}{-0.4}}{bppure}{Get World Transform}{obj/Target: Bocca}{struct/}
\maxx{m1}{s0,tr}
\bpnode{sp}{\dxy{m1}{f}{0.6}{0}}{bpfn}{SpawnActor BP\_Proiettile}{exec/,class/Class,struct/Spawn Transform,obj/Owner: self}{exec/,obj/Return Value}
\bpnode{rl}{\dx{sp}{0.6}{0}}{bpvarh}{Set RiavviaLivello}{exec/,obj/Target: proiettile,bool/RiavviaLivello}{exec/}
\bpnode{d}{\dx{rl}{0.6}{0}}{bpflow}{Delay}{exec/,float/Duration 1.4}{exec/Completed}
\bpnode{s1}{\dx{d}{0.6}{0}}{bpvarh}{Set CadenzaPronta}{exec/,bool/true}{exec/}
\bpexec{f-out1}{s0-in1}\bpexec{s0-out1}{sp-in1}\bpwire{struct}{tr-out1}{sp-in3}\bpexec{sp-out1}{rl-in1}\bpwire{obj}{sp-out2}{rl-in2}\bpexec{rl-out1}{d-in1}\bpexec{d-out1}{s1-in1}
\end{bpgraph}

Nel proiettile, il colpo decide cosa succede:

\begin{bpgraph}
\bpnode{hit}{(0,0)}{bpev}{On Component Hit (Sfera)}{-}{exec/,obj/Other Actor,vec/Normal Impulse}
\bpnode{c}{\dx{hit}{0.6}{0}}{bpcast}{Cast To BP\_Eroe}{exec/,obj/Object}{exec/,exec/Cast Failed,obj/As BP Eroe}
\bpnode[2.4]{br}{\dx{c}{0.6}{0}}{bpflow}{Branch}{exec/,bool/RiavviaLivello}{exec/True,exec/False}
\bpnode{rs}{\dx{br}{0.6}{0}}{bpfn}{UccidiERiavvia}{exec/,float/Ritardo 0.4}{exec/}
\bpnode{dmg}{\sotto{rs}{0}{-0.4}}{bpfn}{Apply Damage}{exec/,obj/Damaged Actor,float/Base Damage 25}{exec/}
\bpnode{des}{\dxy{dmg}{hit}{0.6}{0}}{bpfn}{Destroy Actor}{exec/,obj/self}{exec/}
\bpexec{hit-out1}{c-in1}\bpwire{obj}{hit-out2}{c-in2}\bpexec{c-out1}{br-in1}\bpexec{br-out1}{rs-in1}\bpexec{br-out2}{dmg-in1}\bpexec{rs-out1}{des-in1}\bpexec{dmg-out1}{des-in1}
\bpbox{Il Cast Failed (muro, pavimento) va anch'esso a Destroy Actor}{(hit-nw) (des-se) (dmg-se)}
\end{bpgraph}

\var{UccidiERiavvia} (in \var{BP\_Eroe}): \nodo{Disable Input}, nasconde la mesh, \nodo{Delay} 0,4\,s, \nodo{Open Level (by Name)} con \nodo{Get Current Level Name}. Non toglie una vita: è l'ostacolo ``duro'' del gym, che rimanda all'inizio del livello e basta.

\subsection{La zona che toglie 10 HP al secondo}

La lava del livello 2. Finché si è dentro, ogni secondo arriva un danno di 10 (il primo subito, appena si entra).

\begin{ricetta}
  \item \var{BP\_ZonaDanno}: \nodo{Box Collision} + un piano con materiale emissivo. Variabili: \var{Danno} (10), \var{Intervallo} (1.0), \var{Vittima} (Actor), \var{TimerDanno} (Timer Handle).
  \item In \var{BP\_Eroe}: \var{Vita} (Float, 100) ed \nodo{Event AnyDamage}.
\end{ricetta}

\begin{bpgraph}
\bpnode{ov}{(0,0)}{bpev}{On Component Begin Overlap}{-}{exec/,obj/Other Actor}
\bpnode{s}{\dx{ov}{0.6}{0}}{bpvarh}{Set Vittima}{exec/,obj/}{exec/}
\bpnode{fe}{\dx{s}{0.6}{0}}{bpcustom}{Ferisci}{exec/}{exec/}
\bpnode{tm}{\dx{fe}{0.6}{0}}{bpfn}{Set Timer by Function Name}{exec/,str/Ferisci,float/Time: Intervallo,bool/Looping true}{exec/,struct/Return Value}
\bpnode{st}{\dx{tm}{0.6}{0}}{bpvarh}{Set TimerDanno}{exec/,struct/}{exec/}
\bpexec{ov-out1}{s-in1}\bpwire{obj}{ov-out2}{s-in2}\bpexec{s-out1}{fe-in1}\bpexec{fe-out1}{tm-in1}\bpexec{tm-out1}{st-in1}\bpwire{struct}{tm-out2}{st-in2}
\end{bpgraph}

\begin{bpgraph}
\bpnode{f}{(0,0)}{bpcustom}{Ferisci}{-}{exec/}
\bpnode{ad}{\dx{f}{0.6}{0}}{bpfn}{Apply Damage}{exec/,obj/Damaged Actor: Vittima,float/Base Damage: Danno}{exec/}
\bpnode{end}{($(f-nw|-ad-se)+(0,-0.5)$)}{bpev}{On Component End Overlap}{-}{exec/}
\bpnode{cl}{\dx{end}{0.6}{0}}{bpfn}{Clear and Invalidate Timer by Handle}{exec/,struct/Handle: TimerDanno}{exec/}
\bpexec{f-out1}{ad-in1}\bpexec{end-out1}{cl-in1}
\end{bpgraph}

E nell'eroe:

\begin{bpgraph}
\bpnode{ad}{(0,0)}{bpev}{Event AnyDamage}{-}{exec/,float/Damage}
\bpnode{v}{\dx{ad}{0.6}{0}}{bpvarh}{Set Vita}{exec/,float/Clamp(Vita $-$ Damage ; 0 ; 100)}{exec/,float/}
\bpnode{c}{\sotto{v}{0}{-0.4}}{bppure}{Vita $<=$ 0}{float/}{bool/}
\maxx{m1}{v,c}
\bpnode[2.4]{br}{\dxy{m1}{ad}{0.6}{0}}{bpflow}{Branch}{exec/,bool/Condition}{exec/True,exec/False}
\bpnode{r}{\dx{br}{0.6}{0}}{bpfn}{UccidiERiavvia}{exec/,float/0.5}{exec/}
\bpexec{ad-out1}{v-in1}\bpwire{float}{v-out2}{c-in1}\bpexec{v-out1}{br-in1}\bpwire{bool}{c-out1}{br-in2}\bpexec{br-out1}{r-in1}
\end{bpgraph}

Notate che un \emph{Damage} \emph{negativo} cura: è così che funziona il pickup della salute (capitolo~\ref{sec:pickup}).

\subsection{Zona mortale}

\var{BP\_ZonaMortale}: overlap $\rightarrow$ \nodo{Cast To BP\_Eroe} $\rightarrow$ \var{UccidiERiavvia}(0,3). Serve per le trappole ``istantanee'' (spuntoni); per il vuoto si usa la zona di caduta del capitolo~5, che rispetta le vite.

\subsection{Piattaforma mobile}

Si muove avanti e indietro fra la posizione iniziale e \var{FineOffset} in 2\,s, con 0,5\,s di pausa agli estremi. È il caso d'uso perfetto per una \nodo{Timeline}.

\begin{bpgraph}
\bpnode{bp}{(0,0)}{bpev}{Event BeginPlay}{-}{exec/}
\bpnode{s}{\dx{bp}{0.6}{0}}{bpvarh}{Set Inizio}{exec/,vec/Get Actor Location}{exec/}
\bpnode{tl}{\dx{s}{0.6}{0}}{bptl}{TL\_Va (2 s; curva 0$\rightarrow$1 morbida)}{exec/Play from Start,exec/Reverse from End}{exec/Update,exec/Finished,float/Alpha,enum/Direction}
\bpnode{lerp}{\sotto{tl}{0}{-0.4}}{bppure}{Lerp (Vector)}{vec/A: Inizio,vec/B: Inizio + FineOffset,float/Alpha}{vec/}
\maxx{m1}{tl,lerp}
\bpnode{set}{\dxy{m1}{bp}{0.6}{0}}{bpfn}{Set Actor Location}{exec/,vec/New Location}{exec/}
\bpexec{bp-out1}{s-in1}\bpexec{s-out1}{tl-in1}\bpexec{tl-out1}{set-in1}\bpwire{float}{tl-out3}{lerp-in3}\bpwire{vec}{lerp-out1}{set-in2}
\end{bpgraph}

Su \emph{Finished}: \nodo{Delay} 0,5 $\rightarrow$ \nodo{Branch} su \emph{Direction == Forward} $\rightarrow$ \emph{Reverse from End}, altrimenti \emph{Play from Start}. Così la piattaforma fa avanti e indietro all'infinito. Il personaggio la segue da solo: il Character Movement ``eredita'' la base su cui sta in piedi (\emph{Impart Base Velocity}).

\subsection{Piattaforma che crolla}

Tre stati: \emph{Ferma} $\rightarrow$ (qualcuno ci sale) \emph{Trema} per 0,8\,s $\rightarrow$ \emph{Caduta} (nascosta e senza collisione) $\rightarrow$ dopo 3\,s torna \emph{Ferma}. Uno \emph{switch} su un enum \var{E\_StatoCrollo} nel Tick, oppure, più semplice, una catena di \nodo{Delay}:

\begin{bpgraph}
\bpnode{hit}{(0,0)}{bpev}{On Component Hit (Mesh)}{-}{exec/,obj/Other Actor}
\bpnode{do}{\dx{hit}{0.6}{0}}{bpflow}{Do Once}{exec/,exec/Reset}{exec/Completed}
\bpnode{sh}{\dx{do}{0.6}{0}}{bptl}{TL\_Trema (0,8 s)}{exec/Play from Start}{exec/Update,exec/Finished,float/Rumore}
\bpnode{off}{\dx{sh}{0.6}{0}}{bpfn}{Set Actor Hidden + Collision off}{exec/}{exec/}
\bpnode{d}{\dx{off}{0.6}{0}}{bpflow}{Delay}{exec/,float/3.0}{exec/Completed}
\bpnode{on}{\sotto{off}{0}{-0.5}}{bpfn}{Set Actor Hidden false + Collision on}{exec/}{exec/}
\bpexec{hit-out1}{do-in1}\bpexec{do-out1}{sh-in1}\bpexec{sh-out2}{off-in1}\bpexec{off-out1}{d-in1}
\bpdraw[bpexecw] (d-out1) .. controls +(0.8,0) and +(1.5,0) .. ++(0,-1.3) -- ($(on-nw)+(-0.6,-0.7)$) .. controls +(-0.3,0) and +(-0.6,0) .. (on-in1);
\bpdraw[bpexecw] (on-out1) .. controls +(1.5,0) and +(-1,-1.5) .. (do-in2);
\end{bpgraph}
Durante l'\emph{Update} di \var{TL\_Trema} si sposta la mesh di \nodo{Random Unit Vector}~$\cdot$~4\,cm (\nodo{Set Relative Location}); a \emph{Finished} la si rimette a $(0,0,0)$.

\subsection{Trampolino e barra rotante}

\textbf{Trampolino} (\var{BP\_Trampolino}): overlap $\rightarrow$ cast all'eroe $\rightarrow$ \nodo{Launch Character} con \nodo{Rotate Vector} $\big((300, 0, 1300)$, rotazione dell'actor$\big)$, \emph{XY Override} e \emph{Z Override} veri. Ruotando il trampolino nel livello si cambia la direzione del lancio senza toccare i numeri.

\textbf{Barra rotante}: un \nodo{Rotating Movement Component} (Rotation Rate $(0,0,90)$) fa tutto da solo, senza nodi. Per la spinta: \emph{On Component Hit} della barra $\rightarrow$ \nodo{Launch Character} con \emph{Hit Normal}~$\cdot$~($-$900) + $(0,0,300)$.

\ninja{Timer o Delay? Il \nodo{Delay} è comodo ma non si può fermare. Se l'effetto deve smettere quando si esce da una zona (la lava), usate \nodo{Set Timer} e salvatevi l'\emph{handle}: è l'unico modo per cancellarlo.}

\section{Congegni: interruttori, leve e cose che si attivano}

Nel C++ c'è una coppia di classi base: \texttt{ASwitchBase} (piastra a pressione, leva) e \texttt{AActivable} (porta, luce, piattaforma che sale). L'interruttore non sa \emph{cosa} accende: ha solo una lista di bersagli. In Blueprint si fa uguale, con un'interfaccia e un Event Dispatcher.

\begin{ricetta}
  \item Blueprint Interface \var{BPI\_Attivabile} con le funzioni \var{Attiva} e \var{Disattiva} (nessun parametro).
  \item \var{BP\_Interruttore} (classe padre di piastra e leva): variabile \var{Bersagli} (Array di Actor, \emph{Instance Editable}: si scelgono nel livello con il contagocce), \var{Acceso} (Bool), Event Dispatcher \var{OnCambiato(Acceso)}.
  \item \var{BP\_Porta}, \var{BP\_LuceAttivabile}, \var{BP\_Ascensore}: in \emph{Class Settings} $\rightarrow$ \emph{Implemented Interfaces} aggiungete \var{BPI\_Attivabile}.
\end{ricetta}

\subsection{L'interruttore avvisa i bersagli}

\begin{bpgraph}
\bpnode{f}{(0,0)}{bpcustom}{ImpostaStato}{-}{exec/,bool/Nuovo}
\bpnode{s}{\dx{f}{0.6}{0}}{bpvarh}{Set Acceso}{exec/,bool/Nuovo}{exec/}
\bpnode{call}{\dx{s}{0.6}{0}}{bpcustom}{Call OnCambiato}{exec/,bool/Acceso}{exec/}
\bpnode{fe}{\dx{call}{0.6}{0}}{bpflow}{For Each Loop}{exec/,wild/Array: Bersagli}{exec/Loop Body,wild/Array Element,int/Array Index,exec/Completed}
\bpnode[2.4]{br}{\dx{fe}{0.6}{0}}{bpflow}{Branch}{exec/,bool/Acceso}{exec/True,exec/False}
\bpnode{a}{\dx{br}{0.6}{0.3}}{bpfn}{Attiva (Message)}{exec/,obj/Target}{exec/}
\bpnode{d}{\sotto{a}{0}{-0.3}}{bpfn}{Disattiva (Message)}{exec/,obj/Target}{exec/}
\bpexec{f-out1}{s-in1}\bpwire{bool}{f-out2}{s-in2}\bpexec{s-out1}{call-in1}\bpexec{call-out1}{fe-in1}\bpexec{fe-out1}{br-in1}\bpexec{br-out1}{a-in1}\bpexec{br-out2}{d-in1}
\bpwire{obj}{fe-out2}{a-in2}\bpwire{obj}{fe-out2}{d-in2}
\end{bpgraph}

Il nodo con la busta (\emph{Message}) chiama la funzione dell'interfaccia \emph{se} il bersaglio la implementa, e non fa niente altrimenti: niente cast, niente errori se nella lista finisce un actor sbagliato.

\subsection{La piastra a pressione}

\begin{bpgraph}
\bpnode{ov}{(0,0)}{bpev}{On Component Begin Overlap (Trigger)}{-}{exec/,obj/Other Actor}
\bpnode{c}{\sotto{ov}{0}{-0.4}}{bppure}{Other Actor == Get Player Pawn OR Cast To BP\_Cassa riesce}{-}{bool/}
\maxx{m1}{ov,c}
\bpnode[2.4]{br}{\dxy{m1}{ov}{0.6}{0}}{bpflow}{Branch}{exec/,bool/Condition}{exec/True,exec/False}
\bpnode{st}{\dx{br}{0.6}{0}}{bpcustom}{ImpostaStato}{exec/,bool/Nuovo: true}{exec/}
\bpnode{tl}{\dx{st}{0.6}{0}}{bptl}{TL\_Premi (0,25 s)}{exec/Play,exec/Reverse}{exec/Update,float/Alpha}
\bpexec{ov-out1}{br-in1}\bpwire{bool}{c-out1}{br-in2}\bpexec{br-out1}{st-in1}\bpexec{st-out1}{tl-in1}
\end{bpgraph}

L'\emph{Update} di \var{TL\_Premi} abbassa la piastra di 8\,cm (\nodo{Lerp} fra 0 e $-8$ sulla Z relativa). Anche le casse generate premono la piastra: è un enigma gratis. Se la piastra deve tornare su (\var{RestaPremuta} = false), \emph{End Overlap} $\rightarrow$ \nodo{Get Overlapping Actors} è vuoto? $\rightarrow$ \var{ImpostaStato}(false) e \emph{Reverse}.

\subsection{La leva: interagire con il tasto E}

L'eroe non sa niente delle leve. Quando preme \var{IA\_Interact} lancia una sfera davanti a sé e chiede a quello che colpisce: ``sei interagibile?''.

\begin{bpgraph}
\bpnode{ia}{(0,0)}{bpev}{EnhancedInputAction IA\_Interact}{-}{exec/Triggered,exec/Started}
\bpnode{st}{\sotto{ia}{0}{-0.4}}{bppure}{Camera Loc ; Camera Loc + Forward $\cdot$ 260}{-}{vec/Start,vec/End}
\maxx{m1}{ia,st}
\bpnode{sp}{\dxy{m1}{ia}{0.6}{0}}{bpfn}{Sphere Trace By Channel}{exec/,vec/Start,vec/End,float/Radius 45,enum/Trace Channel: Interact}{exec/,struct/Out Hit,bool/Return Value}
\bpnode{bh}{\sotto{sp}{0}{-0.4}}{bpstruct}{Break Hit Result}{struct/}{obj/Hit Actor}
\bpnode{impl}{\dx{bh}{0.6}{0}}{bppure}{Does Implement Interface}{obj/Test Object,class/BPI\_Interagibile}{bool/}
\maxx{m2}{sp,impl}
\bpnode[2.4]{br}{\dxy{m2}{ia}{0.6}{0}}{bpflow}{Branch}{exec/,bool/Condition}{exec/True,exec/False}
\bpnode{msg}{\dx{br}{0.6}{0}}{bpfn}{Interagisci (Message)}{exec/,obj/Target: Hit Actor,obj/Chi: self}{exec/}
\bpexec{ia-out2}{sp-in1}\bpwire{vec}{st-out1}{sp-in2}\bpwire{vec}{st-out2}{sp-in3}\bpwire{struct}{sp-out2}{bh-in1}\bpwire{obj}{bh-out1}{impl-in1}
\bpexec{sp-out1}{br-in1}\bpwire{bool}{impl-out1}{br-in2}\bpexec{br-out1}{msg-in1}
\end{bpgraph}

Nella leva, l'evento dell'interfaccia:

\begin{bpgraph}
\bpnode{ev}{(0,0)}{bpev}{Event Interagisci}{-}{exec/,obj/Chi}
\bpnode{st}{\dx{ev}{0.6}{0}}{bpcustom}{ImpostaStato}{exec/,bool/NOT Acceso}{exec/}
\bpnode{tl}{\dx{st}{0.6}{0}}{bptl}{TL\_Leva (0,3 s)}{exec/Play,exec/Reverse}{exec/Update,float/Alpha}
\bpnode{r}{\dx{tl}{0.6}{0}}{bpfn}{Set Relative Rotation (Perno)}{exec/,rot/Pitch: Lerp($-$40 ; 40 ; Alpha)}{exec/}
\bpexec{ev-out1}{st-in1}\bpexec{st-out1}{tl-in1}\bpexec{tl-out1}{r-in1}
\end{bpgraph}
(Se \var{Acceso} è vero si fa \emph{Play}, altrimenti \emph{Reverse}: un \nodo{Branch} fra \var{ImpostaStato} e la timeline.) La funzione \var{GetPrompt} restituisce ``Tira la leva'': l'HUD la mostra quando la sfera colpisce qualcosa di interagibile, facendo la stessa traccia nel Tick.

\subsection{La porta}

\begin{bpgraph}
\bpnode{a}{(0,0)}{bpev}{Event Attiva}{-}{exec/}
\bpnode{d}{\sotto{a}{0}{-0.4}}{bpev}{Event Disattiva}{-}{exec/}
\maxx{m1}{a,d}
\bpnode{tl}{\dxy{m1}{a}{0.6}{0}}{bptl}{TL\_Apri (1 s)}{exec/Play,exec/Reverse}{exec/Update,exec/Finished,float/Alpha}
\bpnode{l}{\dx{tl}{0.6}{0}}{bpfn}{Set Relative Location (AntaSx)}{exec/,vec/Lerp((0 0 0) ; (0 $-$120 0) ; Alpha)}{exec/}
\bpnode{r}{\dx{l}{0.6}{0}}{bpfn}{Set Relative Location (AntaDx)}{exec/,vec/Lerp((0 0 0) ; (0 120 0) ; Alpha)}{exec/}
\bpexec{a-out1}{tl-in1}\bpexec{d-out1}{tl-in2}\bpexec{tl-out1}{l-in1}\bpexec{l-out1}{r-in1}
\end{bpgraph}

La porta scorrevole del C++ è questa; per quella a battente si usa \nodo{Set Relative Rotation} sulle cerniere con Yaw fra 0 e 95. La luce attivabile fa \nodo{Set Intensity} con lo stesso Alpha; l'ascensore \nodo{Set Actor Location} con \nodo{Lerp} fra la posizione iniziale e \var{FineOffset} (0, 0, 400).

\ninja{Interfaccia o Event Dispatcher? L'interfaccia serve quando \emph{chi chiama} sceglie a chi parlare (la leva parla ai suoi bersagli). Il dispatcher serve quando \emph{chi ascolta} sceglie di chi fidarsi (un HUD che vuole sapere ``una leva qualsiasi è cambiata'' fa \nodo{Bind Event to OnCambiato}). Nel dubbio: interfaccia.}

\section{Le zone}

Tutte le zone del C++ derivano da \texttt{APlatformerZone}: un \nodo{Box Collision}, un piano colorato per vederla e una scritta. In Blueprint fate lo stesso: un \var{BP\_Zona} padre con l'overlap che filtra solo il giocatore e chiama due eventi vuoti, \var{Entra(Eroe)} ed \var{Esce(Eroe)}. Le zone figlie (\emph{Create Child Blueprint Class}) riempiono solo quei due eventi.

\begin{bpgraph}
\bpnode{ov}{(0,0)}{bpev}{On Component Begin Overlap (Box)}{-}{exec/,obj/Other Actor}
\bpnode{c}{\dx{ov}{0.6}{0}}{bpcast}{Cast To Character}{exec/,obj/Object}{exec/,exec/Cast Failed,obj/As Character}
\bpnode{p}{\dx{c}{0.6}{0}}{bppure}{Is Player Controlled}{obj/Target}{bool/}
\bpnode[2.4]{br}{\dxy{p}{ov}{0.6}{0}}{bpflow}{Branch}{exec/,bool/Condition}{exec/True,exec/False}
\bpnode{e}{\dx{br}{0.6}{0}}{bpcustom}{Entra}{exec/,obj/Eroe}{exec/}
\bpexec{ov-out1}{c-in1}\bpwire{obj}{ov-out2}{c-in2}\bpwire{obj}{c-out3}{p-in1}\bpexec{c-out1}{br-in1}\bpwire{bool}{p-out1}{br-in2}\bpexec{br-out1}{e-in1}
\bpbox{BP\_Zona (padre): filtra il giocatore. Uguale per End Overlap $\rightarrow$ Esce}{(ov-nw) (e-se) (c-se)}
\end{bpgraph}

Nel figlio, tasto destro su \var{Entra} $\rightarrow$ \emph{Add call to parent function} se volete anche il comportamento del padre.

\subsection{Zona che modifica il movimento (ghiaccio, fango, bassa gravità)}

All'ingresso si salvano i valori originali, si moltiplicano, e all'uscita (dopo 3\,s, \var{DurataDopoUscita}) si rimettono. Variabili: i moltiplicatori (\var{Velocita}, \var{Salto}, \var{Gravita}, \var{Attrito}, \var{Frenata}) e i salvati.

\begin{bpgraph}
\bpnode{e}{(0,0)}{bpev}{Event Entra}{-}{exec/,obj/Eroe}
\bpnode{mv}{\sotto{e}{0}{-0.4}}{bppure}{Get Character Movement}{obj/Eroe}{obj/Move}
\maxx{m1}{e,mv}
\bpnode{s1}{\dxy{m1}{e}{0.6}{0}}{bpvarh}{Set SalvataVel}{exec/,float/Move.Max Walk Speed}{exec/}
\bpnode{s2}{\dx{s1}{0.6}{0}}{bpvarh}{Set SalvatoAttrito}{exec/,float/Move.Ground Friction}{exec/}
\bpnode{m1}{\dx{s2}{0.6}{0}}{bpvarh}{Set Max Walk Speed}{exec/,obj/Move,float/SalvataVel $\cdot$ Velocita}{exec/}
\bpnode{m2}{\dx{m1}{0.6}{0}}{bpvarh}{Set Ground Friction}{exec/,obj/Move,float/SalvatoAttrito $\cdot$ Attrito}{exec/}
\bpexec{e-out1}{s1-in1}\bpexec{s1-out1}{s2-in1}\bpexec{s2-out1}{m1-in1}\bpexec{m1-out1}{m2-in1}\bpwire{obj}{e-out2}{mv-in1}
\end{bpgraph}

(Lo stesso per Jump Z Velocity, Gravity Scale e Braking Deceleration Walking.) Il ghiaccio è \var{Attrito} 0,05 e \var{Frenata} 0,1; il fango \var{Velocita} 0,5 e \var{Salto} 0,7; la luna \var{Gravita} 0,3.

\ninja{Salvate i valori \emph{all'ingresso}, non nel BeginPlay della zona: se due zone si sovrappongono, o un'abilità ha cambiato i valori, ripristinate quello che c'era davvero.}

\subsection{Teletrasporto}

\var{Entra} $\rightarrow$ \nodo{Teleport} (Target: Eroe, Dest Location: \var{Destinazione}, Dest Rotation: rotazione dell'eroe) $\rightarrow$ se \var{FermaVelocita}, \nodo{Set Velocity} (0,0,0) sul Movement. \var{Destinazione} è un Vector con \emph{Show 3D Widget} acceso: nel livello compare un rombo da trascinare (attenzione, è relativo all'actor: convertitelo con \nodo{Transform Location}).

\subsection{Il negozio: spendere monete}

\begin{bpgraph}
\bpnode{e}{(0,0)}{bpev}{Event Entra}{-}{exec/,obj/Eroe}
\bpnode{c}{\sotto{e}{0}{-0.4}}{bppure}{GI.Monete $>=$ Costo AND NOT GI.HaAbilita(Premio)}{-}{bool/}
\maxx{m1}{e,c}
\bpnode[2.4]{br}{\dxy{m1}{e}{0.6}{0}}{bpflow}{Branch}{exec/,bool/Condition}{exec/True,exec/False}
\bpnode{pay}{\dx{br}{0.6}{0}}{bpvarh}{Set GI.Monete}{exec/,int/Monete $-$ Costo}{exec/}
\bpnode{ab}{\dx{pay}{0.6}{0}}{bpfn}{SbloccaAbilita}{exec/,obj/GI,enum/Premio}{exec/}
\bpnode{no}{\sotto{pay}{0}{-0.4}}{bpfn}{Mostra Messaggio}{exec/,str/"Servono 5 monete"}{exec/}
\bpexec{e-out1}{br-in1}\bpwire{bool}{c-out1}{br-in2}\bpexec{br-out1}{pay-in1}\bpexec{pay-out1}{ab-in1}\bpexec{br-out2}{no-in1}
\end{bpgraph}

\subsection{Traguardo e timer del livello}

\textbf{Traguardo} (\var{BP\_Traguardo}): \var{Entra} $\rightarrow$ \nodo{Do Once} $\rightarrow$ messaggio ``Livello completato!'' $\rightarrow$ salva il record (\nodo{Find} nella Map \var{Record} con il nome del livello; se il tempo è minore, \nodo{Add} lo sostituisce) $\rightarrow$ \nodo{Delay} 1\,s $\rightarrow$ \var{VaiAlLivello(ProssimoLivello)}.

\textbf{Timer} (\var{BP\_TimerLivello}, livello 3: 150\,s). Un \nodo{Text Render} grande sospeso nel livello, e un timer da un decimo di secondo:

\begin{bpgraph}
\bpnode{bp}{(0,0)}{bpev}{Event BeginPlay}{-}{exec/}
\bpnode{s}{\dx{bp}{0.6}{0}}{bpvarh}{Set TempoRimasto}{exec/,float/TempoLimite (150)}{exec/}
\bpnode{tm}{\dx{s}{0.6}{0}}{bpfn}{Set Timer by Event}{exec/,del/Event,float/Time 0.1,bool/Looping true}{exec/}
\bpnode{ev}{\sotto{bp}{0}{-1.6}}{bpcustom}{Custom Event: Scorri}{-}{exec/,del/}
\bpnode{dec}{\dx{ev}{0.6}{0}}{bpvarh}{Set TempoRimasto}{exec/,float/TempoRimasto $-$ 0.1}{exec/,float/}
\bpnode{txt}{\dx{dec}{0.6}{0}}{bpfn}{Set Text (Text Render)}{exec/,text/Format: "\{0\} s" (1 decimale)}{exec/}
\bpnode{col}{\dx{txt}{0.6}{0}}{bpfn}{Set Text Render Color}{exec/,struct/Select: rosso se $<$ 10 s}{exec/}
\bpexec{bp-out1}{s-in1}\bpexec{s-out1}{tm-in1}\bpdraw[bpdataw, draw=bpc-del] (ev-out2) .. controls +(2,0.5) and +(-1,0) .. (tm-in2);
\bpexec{ev-out1}{dec-in1}\bpexec{dec-out1}{txt-in1}\bpexec{txt-out1}{col-in1}
\end{bpgraph}
Dopo il colore, un \nodo{Branch} su \var{TempoRimasto} $<=$ 0: \nodo{Clear Timer}, messaggio ``Tempo scaduto'' e \var{UccidiERiavvia}. Il \nodo{Format Text} con \{0\} si collega a un \nodo{To Text (Float)} con \emph{Maximum Fractional Digits} = 1.

\section{I pickup e i potenziamenti}\label{sec:pickup}

Un solo padre, \var{BP\_Pickup}, fa girare e ondeggiare l'oggetto e gestisce la raccolta; i figli decidono soltanto cosa succede quando lo si prende. Nel C++: \texttt{APlatformerPickup} con \texttt{ACoinPickup}, \texttt{AAbilityPickup}, \texttt{AHealthPickup}, \texttt{ABatteryPickup}.

\begin{ricetta}
  \item \var{BP\_Pickup}: \var{Sfera} (Sphere Collision, raggio 60, root), \var{Mesh}, \var{Luce} (Point Light, intensità 8, raggio piccolo), \var{Scritta} (Text Render, rivolta alla telecamera). Variabili: \var{Giro} (120\textdegree/s), \var{Onda} (12\,cm), \var{VelOnda} (0,8), \var{Orologio}, \var{ZBase}.
  \item Funzione vuota \var{Raccolto(Eroe)} che ritorna un Bool: \emph{True} se il pickup va consumato.
\end{ricetta}

\subsection{Girare e ondeggiare}

\begin{bpgraph}
\bpnode{t}{(0,0)}{bpev}{Event Tick}{-}{exec/,float/Delta Seconds}
\bpnode{o}{\dx{t}{0.6}{0}}{bpvarh}{Set Orologio}{exec/,float/Orologio + dt}{exec/}
\bpnode{r}{\dx{o}{0.6}{0}}{bpfn}{Add Local Rotation}{exec/,obj/Mesh,rot/(0 ; Giro $\cdot$ dt ; 0)}{exec/}
\bpnode{z}{\sotto{r}{0}{-0.4}}{bppure}{ZBase + Sin(Orologio $\cdot$ VelOnda $\cdot$ 2$\pi$) $\cdot$ Onda}{-}{float/}
\maxx{m1}{r,z}
\bpnode{l}{\dxy{m1}{t}{0.6}{0}}{bpfn}{Set Relative Location}{exec/,obj/Mesh,vec/(0 0 Z)}{exec/}
\bpexec{t-out1}{o-in1}\bpexec{o-out1}{r-in1}\bpexec{r-out1}{l-in1}\bpwire{float}{z-out1}{l-in3}
\end{bpgraph}

Il \nodo{Sin} lavora in radianti: per questo il $2\pi$. Con \var{VelOnda} 0,8 l'oggetto fa quasi un'onda al secondo.

\subsection{La raccolta (nel padre)}

\begin{bpgraph}
\bpnode{ov}{(0,0)}{bpev}{On Component Begin Overlap (Sfera)}{-}{exec/,obj/Other Actor}
\bpnode{c}{\dx{ov}{0.6}{0}}{bpcast}{Cast To BP\_Eroe}{exec/,obj/Object}{exec/,exec/Cast Failed,obj/As BP Eroe}
\bpnode{rc}{\dx{c}{0.6}{0}}{bpfn}{Raccolto}{exec/,obj/Eroe}{exec/,bool/Consuma}
\bpnode[2.4]{br}{\dx{rc}{0.6}{0}}{bpflow}{Branch}{exec/,bool/Condition}{exec/True,exec/False}
\bpnode{snd}{\dx{br}{0.6}{0}}{bpfn}{Play Sound at Location}{exec/,obj/Sound,vec/Location}{exec/}
\bpnode{d}{\dx{snd}{0.6}{0}}{bpfn}{Destroy Actor}{exec/}{exec/}
\bpexec{ov-out1}{c-in1}\bpwire{obj}{ov-out2}{c-in2}\bpexec{c-out1}{rc-in1}\bpwire{obj}{c-out3}{rc-in2}\bpexec{rc-out1}{br-in1}\bpwire{bool}{rc-out2}{br-in2}\bpexec{br-out1}{snd-in1}\bpexec{snd-out1}{d-in1}
\end{bpgraph}

\subsection{I figli: ridefinire \var{Raccolto}}

Nel figlio, nel pannello \emph{My Blueprint} $\rightarrow$ \emph{Functions} $\rightarrow$ \emph{Override} $\rightarrow$ \var{Raccolto}.

\begin{bpgraph}
\bpnode{f}{(0,0)}{bpcustom}{Raccolto (Moneta)}{-}{exec/,obj/Eroe}
\bpnode{m}{\dx{f}{0.6}{0}}{bpvarh}{Set GI.Monete}{exec/,int/Monete + Valore}{exec/}
\bpnode{r}{\dx{m}{0.6}{0}}{bpflow}{Return Node}{exec/,bool/Consuma: true}{-}
\bpnode{f2}{\sotto{f}{0}{-0.5}}{bpcustom}{Raccolto (Abilità)}{-}{exec/,obj/Eroe}
\bpnode{a}{\dx{f2}{0.6}{0}}{bpfn}{SbloccaAbilita}{exec/,obj/GI,enum/Abilita}{exec/}
\bpnode{msg}{\dx{a}{0.6}{0}}{bpfn}{Mostra Messaggio}{exec/,text/Descrizione}{exec/}
\bpnode{r2}{\dx{msg}{0.6}{0}}{bpflow}{Return Node}{exec/,bool/Consuma: true}{-}
\bpexec{f-out1}{m-in1}\bpexec{m-out1}{r-in1}\bpexec{f2-out1}{a-in1}\bpexec{a-out1}{msg-in1}\bpexec{msg-out1}{r2-in1}
\end{bpgraph}

\begin{bpgraph}
\bpnode{f}{(0,0)}{bpcustom}{Raccolto (Salute)}{-}{exec/,obj/Eroe}
\bpnode{c}{\sotto{f}{0}{-0.4}}{bppure}{Eroe.Vita $<$ 100}{-}{bool/}
\maxx{m1}{f,c}
\bpnode[2.4]{br}{\dxy{m1}{f}{0.6}{0}}{bpflow}{Branch}{exec/,bool/Condition}{exec/True,exec/False}
\bpnode{d}{\dx{br}{0.6}{0}}{bpfn}{Apply Damage}{exec/,obj/Damaged Actor: Eroe,float/Base Damage: $-$40}{exec/}
\bpnode{r}{\dx{d}{0.6}{0}}{bpflow}{Return Node}{exec/,bool/Consuma: true}{-}
\bpnode{rf}{\sotto{d}{0}{-0.4}}{bpflow}{Return Node}{exec/,bool/Consuma: false}{-}
\bpexec{f-out1}{br-in1}\bpwire{bool}{c-out1}{br-in2}\bpexec{br-out1}{d-in1}\bpexec{d-out1}{r-in1}\bpexec{br-out2}{rf-in1}
\end{bpgraph}

Con la vita piena la salute \emph{non} si consuma: resta lì per dopo. La batteria fa lo stesso con \var{Batteria} della torcia (+15\,s, \nodo{Min} con il massimo).

\subsection{Il potenziamento con il testo}

Il brief chiede potenziamenti ``semplici, con un testo''. Il \var{BP\_PickupAbilita} ha una variabile \var{Abilita} (\var{E\_Abilita}, \emph{Instance Editable}) e nel \nodo{Construction Script} scrive da solo il proprio cartello:

\begin{bpgraph}
\bpnode{cs}{(0,0)}{bpev}{Construction Script}{-}{exec/}
\bpnode{sw}{\sotto{cs}{0}{-0.4}}{bppure}{Select (E\_Abilita)}{enum/Index: Abilita,text/DoppioSalto: "Doppio salto",text/WallRun: "Corsa sul muro",text/Jetpack: "Jetpack (tieni Salto)",text/GeneratoreCasse: "Casse (F)"}{text/}
\maxx{m1}{cs,sw}
\bpnode{t}{\dxy{m1}{cs}{0.6}{0}}{bpfn}{Set Text}{exec/,obj/Target: Scritta,text/Value}{exec/}
\bpnode{col}{\dx{t}{0.6}{0}}{bpfn}{Set Vector Parameter Value on Materials}{exec/,obj/Mesh,name/"Colore",vec/colore dell'abilità}{exec/}
\bpexec{cs-out1}{t-in1}\bpwire{text}{sw-out1}{t-in3}\bpexec{t-out1}{col-in1}
\end{bpgraph}

Il \nodo{Construction Script} gira anche nell'editor: cambiate l'abilità nel pannello Details e il cartello si aggiorna mentre costruite il livello.

\ninja{Un pickup che si ``ricorda'' di essere già stato preso fra un livello e l'altro? Dategli un \var{ID} (Name) e, in BeginPlay, se \var{GI.Raccolti} contiene l'\var{ID}, \nodo{Destroy Actor}. È così che le abilità non ricompaiono dopo averle prese.}

\section{L'HUD: il giocatore deve vedere tutto}

Nel C++ l'HUD è disegnato a mano (\texttt{APlatformerHUD}). In Blueprint si usa un \nodo{Widget Blueprint}, molto più comodo: si disegna nel Designer e si collegano i valori.

\begin{ricetta}
  \item \var{WBP\_HUD}: in alto a sinistra un \nodo{Vertical Box} con testi per vite (un'icona ripetuta tre volte), monete e timer; in basso tre \nodo{Progress Bar} (vita, carburante, batteria) e un testo per le casse; al centro un mirino (un \nodo{Image} piccolo) visibile solo in prima persona; un \nodo{Text Block} \var{Messaggio} per gli avvisi.
  \item In \var{BP\_Eroe}, \nodo{Event BeginPlay} $\rightarrow$ \nodo{Create Widget} (\var{WBP\_HUD}, Owning Player: \nodo{Get Player Controller}) $\rightarrow$ \nodo{Add to Viewport}, e salvate il widget in \var{HUD}.
\end{ricetta}

\subsection{Le barre: binding}

Selezionate la barra del carburante, \emph{Percent} $\rightarrow$ \emph{Bind} $\rightarrow$ \emph{Create Binding}. Si apre una funzione pura che UE chiama a ogni frame:

\begin{bpgraph}
\bpnode{f}{(0,0)}{bpcustom}{Get Percent Carburante}{-}{exec/}
\bpnode{p}{\sotto{f}{0}{-0.4}}{bppure}{Get Owning Player Pawn}{-}{obj/}
\bpnode{c}{\dxy{p}{f}{0.6}{0}}{bpcast}{Cast To BP\_Eroe}{exec/,obj/Object}{exec/,exec/Cast Failed,obj/As BP Eroe}
\bpnode{d}{\sotto{c}{0}{-0.4}}{bppure}{Carburante / CarburanteMax}{obj/Eroe}{float/}
\maxx{m1}{c,d}
\bpnode{r}{\dxy{m1}{f}{0.6}{0}}{bpflow}{Return Node}{exec/,float/Return Value}{-}
\bpexec{f-out1}{c-in1}\bpwire{obj}{p-out1}{c-in2}\bpwire{obj}{c-out3}{d-in1}\bpexec{c-out1}{r-in1}\bpwire{float}{d-out1}{r-in2}
\end{bpgraph}

Allo stesso modo la \emph{Visibility} della barra: \emph{Visible} se l'eroe ha il jetpack, \emph{Collapsed} altrimenti (un \nodo{Select} sull'enum \emph{ESlateVisibility}). Il testo delle casse: \nodo{Format Text} ``Casse: \{0\}/3'' oppure ``Ricarica\dots\ \{1\} s''.

\subsection{Gli avvisi: eventi invece del binding}

Il binding è comodo ma gira a ogni frame. Per le cose che cambiano di rado (un'abilità sbloccata, un messaggio) è meglio un evento:

\begin{bpgraph}
\bpnode{ev}{(0,0)}{bpcustom}{Custom Event: MostraMessaggio}{-}{exec/,text/Testo,float/Durata}
\bpnode{t}{\dx{ev}{0.6}{0}}{bpfn}{Set Text}{exec/,obj/Messaggio,text/Testo}{exec/}
\bpnode{v}{\dx{t}{0.6}{0}}{bpfn}{Set Visibility}{exec/,obj/Messaggio,enum/Visible}{exec/}
\bpnode{an}{($(ev-nw|-v-se)+(1.2,-0.5)$)}{bpfn}{Play Animation}{exec/,obj/Anim: Comparsa}{exec/}
\bpnode{rt}{\dx{an}{0.6}{0}}{bpflow}{Retriggerable Delay}{exec/,float/Durata}{exec/Completed}
\bpnode{h}{\dx{rt}{0.6}{0}}{bpfn}{Set Visibility}{exec/,obj/Messaggio,enum/Hidden}{exec/}
\bpexec{ev-out1}{t-in1}\bpexec{t-out1}{v-in1}\bpdraw[bpexecw] (v-out1) .. controls +(1.2,0) and +(1,0) .. ($(v-se)+(0.3,-0.25)$) -- ($(an-nw)+(-0.5,0.25)$) .. controls +(-1,0) and +(-1,0) .. (an-in1);\bpexec{an-out1}{rt-in1}\bpexec{rt-out1}{h-in1}
\end{bpgraph}

Il \nodo{Retriggerable Delay} riparte da capo se arriva un secondo messaggio prima che il primo sparisca: niente messaggi che si spengono troppo presto. Nel \nodo{Event Construct} del widget fate \nodo{Bind Event to OnAbilitaSbloccata} della Game Instance e chiamate \var{MostraMessaggio} da lì.

\ninja{Il widget non deve \emph{mai} modificare il gioco: legge e mostra. Se un pulsante dell'HUD deve fare qualcosa, chiama una funzione dell'eroe o della Game Instance. Così potete cambiare tutto l'HUD senza rompere niente.}

\section{L'ambiente che si muove}

Piccoli actor di contorno, ma ognuno insegna un nodo.

\subsection{Il mulino: rotazione continua}

\var{BP\_Mulino}: torre, \var{Mozzo} (Scene) con quattro pale figlie. Nel Tick: \nodo{Add Local Rotation} sul \var{Mozzo} con Roll = 60~$\cdot$~dt. In più, la testa del mulino oscilla lentamente di $\pm 60$\textdegree:

\begin{bpgraph}
\bpnode{t}{(0,0)}{bpev}{Event Tick}{-}{exec/,float/Delta Seconds}
\bpnode{p}{\dx{t}{0.6}{0}}{bpfn}{Add Local Rotation}{exec/,obj/Mozzo,rot/Roll: 60 $\cdot$ dt}{exec/}
\bpnode{y}{\sotto{p}{0}{-0.4}}{bppure}{Sin(Get Game Time $\cdot$ 0.2) $\cdot$ 60}{-}{float/Yaw}
\maxx{m1}{p,y}
\bpnode{r}{\dxy{m1}{t}{0.6}{0}}{bpfn}{Set Relative Rotation}{exec/,obj/Testa,rot/(0 ; Yaw ; 0)}{exec/}
\bpexec{t-out1}{p-in1}\bpexec{p-out1}{r-in1}\bpwire{float}{y-out1}{r-in3}
\end{bpgraph}

\subsection{La luce che tremola (fuoco, candela)}

\begin{bpgraph}
\bpnode{t}{(0,0)}{bpev}{Event Tick}{-}{exec/,float/Delta Seconds}
\bpnode{rn}{\sotto{t}{0}{-0.4}}{bppure}{Random Float in Range}{float/Min 0.7,float/Max 1.15}{float/}
\bpnode{fi}{\dx{rn}{0.6}{0}}{bppure}{FInterp To}{float/Current: Fattore,float/Target,float/Delta Time,float/Interp Speed 12}{float/}
\bpnode{s}{\dxy{fi}{t}{0.6}{0}}{bpvarh}{Set Fattore}{exec/,float/}{exec/,float/}
\bpnode{i}{\dx{s}{0.6}{0}}{bpfn}{Set Intensity}{exec/,obj/Luce,float/BaseIntensity $\cdot$ Fattore}{exec/}
\bpexec{t-out1}{s-in1}\bpwire{float}{rn-out1}{fi-in2}\bpwire{float}{fi-out1}{s-in2}\bpexec{s-out1}{i-in1}
\end{bpgraph}

Un numero a caso \emph{diretto} farebbe una luce stroboscopica. Il \nodo{FInterp To} verso un bersaglio casuale è il trucco: il fuoco ``respira''.

\subsection{La staccionata lungo una spline}

\var{BP\_Staccionata}: un \nodo{Spline Component} e un \nodo{Instanced Static Mesh} per i pali. Tutto nel \nodo{Construction Script}, così la staccionata si ridisegna mentre trascinate i punti della spline.

\begin{bpgraph}
\bpnode{cs}{(0,0)}{bpev}{Construction Script}{-}{exec/}
\bpnode{cl}{\dx{cs}{0.6}{0}}{bpfn}{Clear Instances}{exec/,obj/Pali}{exec/}
\bpnode{n}{\sotto{cl}{0}{-0.4}}{bppure}{Truncate(Get Spline Length / 200)}{-}{int/N}
\maxx{m1}{cl,n}
\bpnode{fl}{\dxy{m1}{cs}{0.6}{0}}{bpflow}{For Loop}{exec/,int/First 0,int/Last: N}{exec/Loop Body,int/Index,exec/Completed}
\bpnode{tr}{\sotto{fl}{0}{-0.4}}{bppure}{Get Transform at Distance Along Spline}{float/Index $\cdot$ 200,enum/Local}{struct/}
\maxx{m2}{fl,tr}
\bpnode{ad}{\dxy{m2}{cs}{0.6}{0}}{bpfn}{Add Instance}{exec/,obj/Pali,struct/Instance Transform}{exec/}
\bpexec{cs-out1}{cl-in1}\bpexec{cl-out1}{fl-in1}\bpwire{int}{n-out1}{fl-in3}\bpexec{fl-out1}{ad-in1}\bpwire{int}{fl-out2}{tr-in1}\bpwire{struct}{tr-out1}{ad-in3}
\end{bpgraph}

Le traverse (\var{Rails} = 2) si fanno con un \nodo{Spline Mesh Component} per ogni tratto fra due pali: \nodo{Add Spline Mesh Component} $\rightarrow$ \nodo{Set Start and End} con posizioni e tangenti lette dalla spline.

\subsection{Lo spargitore: sassi e cespugli a caso}

\var{BP\_Spargi}: un \nodo{Box} che delimita l'area e un \nodo{Instanced Static Mesh}. Nel \nodo{Construction Script}, un \nodo{For Loop} da 0 a \var{Quanti} (30): punto a caso nel box (\nodo{Random Point in Bounding Box From Stream}, con un \nodo{Random Stream} di seme 7 così il risultato non cambia a ogni modifica), un \nodo{Line Trace} verso il basso per appoggiarlo al terreno, scala a caso fra 0,8 e 1,3, yaw a caso, \nodo{Add Instance}.

\ninja{\nodo{Instanced Static Mesh}: cento sassi uguali disegnati con una sola chiamata alla scheda video. Cento actor separati sono cento chiamate. Per il ``contorno'' è sempre la scelta giusta.}

\section{Le mosse da ninja}

Le regole che separano un Blueprint che ``funziona'' da uno che si può ancora leggere fra un mese.

\begin{enumerate}
  \item \textbf{Pure e impure.} Un nodo verde (puro) viene ricalcolato \emph{ogni volta} che un pin lo legge. Se \nodo{Random Float} è puro e lo collegate a due nodi, avete due numeri diversi. Se un valore deve essere lo stesso, salvatelo in una variabile locale.
  \item \textbf{Is Valid prima di usare un riferimento.} Actor distrutti (casse scadute, proiettili) lasciano riferimenti vuoti. Il nodo \nodo{Is Valid} con i due pin di esecuzione è più leggibile del \nodo{Branch}.
  \item \textbf{Interfacce invece di cast a catena.} Se vi ritrovate con \nodo{Cast To A}, se fallisce \nodo{Cast To B}, se fallisce \nodo{Cast To C}\dots\ volete un'interfaccia (capitolo 9).
  \item \textbf{Timer invece del Tick.} Il Tick gira 60--120 volte al secondo. Il posto sicuro, il cambio modello, la torretta che decide se sparare: 5 volte al secondo bastano. Il Tick tenetelo per ciò che si \emph{vede} muoversi.
  \item \textbf{Delay non si ferma.} Usate \nodo{Set Timer} con l'handle quando l'effetto può essere interrotto.
  \item \textbf{Un componente per abilità.} \var{AC\_Movimento}, \var{AC\_Torcia}, \var{AC\_PostoSicuro}, \var{AC\_Casse}: si attaccano e staccano, si provano da soli, e il Character resta corto. È la stessa scelta fatta nel C++.
  \item \textbf{Funzioni e Collapse.} Selezionate un gruppo di nodi $\rightarrow$ tasto destro $\rightarrow$ \emph{Collapse to Function}. Un grafo che non sta in uno schermo va spezzato.
  \item \textbf{Commenti (tasto C).} Ogni gruppo di nodi con un titolo, come i riquadri beige di questa guida.
  \item \textbf{Reroute node.} Doppio clic su un filo crea un punto di passaggio: niente fili che attraversano i nodi.
  \item \textbf{Instance Editable e Expose on Spawn.} Le variabili che il level designer deve toccare (danno, tempi, bersagli) vanno esposte con l'occhio aperto. I numeri magici dentro i nodi sono invisibili a chi costruisce il livello.
\end{enumerate}

\subsection{Debug}

\begin{itemize}
  \item \nodo{Print String} con \emph{Key} fissa (per esempio ``carburante''): il messaggio si aggiorna al suo posto invece di riempire lo schermo.
  \item Durante il Play, nel Blueprint scegliete l'istanza da osservare in alto (\emph{Debug Object}): i fili attivi si illuminano e passando sopra un pin si legge il valore.
  \item \emph{F9} su un nodo mette un breakpoint. \emph{F10} va avanti di un nodo.
  \item Le tracce hanno \emph{Draw Debug Type}: \emph{For Duration} mostra i raggi del wall run e del posto sicuro. È il modo più veloce di capire perché un muro ``non prende''.
  \item Console (tasto \`{}): \texttt{slomo 0.2} rallenta il gioco, \texttt{stat fps} mostra i frame, \texttt{show collision} mostra tutte le collisioni.
\end{itemize}

\subsection{Corrispondenze C++ $\leftrightarrow$ Blueprint}

\begin{center}
\small
\begin{tabularx}{\linewidth}{>{\raggedright\arraybackslash}X>{\raggedright\arraybackslash}p{6.2cm}}
\toprule
Nel codice & In Blueprint \\
\midrule
\texttt{UPlatformerRunSubsystem} & \var{GI\_Platform} (Game Instance) \\
\texttt{UPlatformerAbilitiesComponent} & \var{AC\_Movimento} (cap.~3) \\
\texttt{UPlatformerJetpackComponent} & Tick del jetpack in \var{AC\_Movimento} (cap.~3.6) \\
\texttt{UPlatformerTorchComponent} & \var{AC\_Torcia} (cap.~4) \\
\texttt{UPlatformerSafeSpotComponent}, \texttt{AFallZone} & \var{AC\_PostoSicuro}, \var{BP\_ZonaCaduta} (cap.~5) \\
\texttt{UPlatformerCrateSpawnerComponent}, \texttt{ASpawnedCrate} & \var{AC\_Casse}, \var{BP\_Cassa} (cap.~6) \\
\texttt{AScaleByDistanceActor}, \texttt{AModelCyclerActor}, \texttt{ARandomPresenceActor} & cap.~7 \\
\texttt{ATurret}, \texttt{APlatformerProjectile}, \texttt{ADamageZone}, \texttt{AKillZone} & cap.~8 \\
\texttt{AMovingPlatform}, \texttt{ACrumblingPlatform}, \texttt{AJumpPad}, \texttt{ASpinningObstacle} & cap.~8 \\
\texttt{ASwitchBase}, \texttt{AActivable} e figli & \var{BP\_Interruttore}, \var{BPI\_Attivabile} (cap.~9) \\
\texttt{AModifierZone}, \texttt{ATeleportZone}, \texttt{AShopZone}, \texttt{ALevelGoal}, \texttt{ALevelTimer} & cap.~10 \\
\texttt{APlatformerPickup} e figli & \var{BP\_Pickup} e figli (cap.~11) \\
\texttt{APlatformerHUD} & \var{WBP\_HUD} (cap.~12) \\
\texttt{AWindmill}, \texttt{AFlickeringLight}, \texttt{ASplineFence}, \texttt{AScatterArea} & cap.~13 \\
\texttt{APlatformerGameMode} & \var{GM\_Platform}: Default Pawn = \var{BP\_Eroe} \\
\bottomrule
\end{tabularx}
\end{center}

\ninja{L'ultima mossa: rifate una meccanica \emph{senza} guardare la guida. Se vi bloccate, è il punto esatto che non avevate capito; tornate solo a quella pagina.}


\end{document}
